\documentclass[11pt,a4paper]{article}
\usepackage[margin=23mm]{geometry}
\usepackage[T1]{fontenc}
\usepackage{lmodern,amsmath,amssymb,amsthm,booktabs,array,longtable,fancyhdr}
\usepackage[colorlinks=true,linkcolor=blue,citecolor=blue,urlcolor=blue]{hyperref}
\hypersetup{pdftitle={Cubic twists of Legendre moments: cohomological descent and Frobenius trace formulas},pdfauthor={Rushikesh Murlidhar Walke},pdfsubject={Research preprint on cubic Legendre moments, local factors and conductors}}
\setlength{\parindent}{1.2em}\setlength{\parskip}{3pt}
\setlength{\emergencystretch}{3em}\setlength{\headheight}{14pt}
\pagestyle{fancy}\fancyhf{}
\fancyhead[L]{\small Cubic twists of Legendre moments}
\fancyhead[R]{\small Research preprint}\fancyfoot[C]{\thepage}
\newcommand{\Q}{\mathbb Q}\newcommand{\F}{\mathbb F}\newcommand{\Z}{\mathbb Z}
\newcommand{\C}{\mathbb C}\newcommand{\R}{\mathbb R}\newcommand{\PP}{\mathbb P}
\newcommand{\Sym}{\operatorname{Sym}}\newcommand{\Tr}{\operatorname{Tr}}
\newcommand{\Frob}{\operatorname{Frob}}\newcommand{\Gal}{\operatorname{Gal}}
\newcommand{\nt}{\mathrm{nt}}
\newcommand{\Sw}{\operatorname{Sw}}\newcommand{\HT}{\operatorname{HT}}
\newcommand{\GL}{\operatorname{GL}}\newcommand{\CH}{\operatorname{CH}}
\newcommand{\sgn}{\operatorname{sgn}}\newcommand{\im}{\operatorname{im}}
\newcommand{\et}{\mathrm{\acute et}}\newcommand{\Hpar}{H^1_!}
\newcommand{\leg}[2]{\left(\frac{#1}{#2}\right)}
\newtheorem{theorem}{Theorem}[section]
\newtheorem{lemma}[theorem]{Lemma}\newtheorem{proposition}[theorem]{Proposition}
\newtheorem{corollary}[theorem]{Corollary}
\theoremstyle{remark}\newtheorem{remark}[theorem]{Remark}
\newtheorem{observation}[theorem]{Numerical observation}
\newtheorem{conjecture}[theorem]{Conjecture}
\title{Cubic twists of Legendre moments:\\cohomological descent and Frobenius trace formulas}
\author{Rushikesh Murlidhar Walke\\\small Independent Researcher, Maharashtra, India\\\small Corresponding author: \href{mailto:Rushikesh@rmwalke.com}{\texttt{Rushikesh@rmwalke.com}}}
\date{Research preprint --- 30 October 2026}
\begin{document}
\maketitle
\begin{abstract}
We construct the parabolic cohomology representations attached to cubic
twists of the second and fourth moments of the Legendre family. For the
fourth moment, an explicit arithmetic involution gives two rational
four-dimensional descents related by the quadratic character of
$\Q(\sqrt{-3})$. We prove their alternating pairing and derive all
good-prime Euler factors. Deligne's weight theorem gives good-prime
purity of weight five. An explicit auxiliary-level construction realizes
both descents by rational Chow projectors on a smooth projective fivefold;
their Hodge numbers are $h^{5,0}=h^{0,5}=2$, with no other Hodge types.
At inert primes a norm-one-torus calculation
determines the trace, the boundary contribution and the sign matching
the earlier finite-sum recipe. The cubic second-moment modular comparison
uses explicit finite-check hypotheses and the supplied bounded exact
certificate for the displayed rows and Sturm comparison.
We compute inertia invariants and Swan conductors at both bad primes.
The literal plus descent has conductor $432$, and its quadratic twist has
conductor $1296$. A local-monodromy and Hodge--Tate argument proves strong
absolute irreducibility for coefficient primes $\ell\ne3$.
We retain the corrected hypergeometric normalization and distinguish these
written local proofs from the remaining numerical special-value observations.
The underlying cover is compared explicitly with existing literature;
no general BSD consequence or literature-priority certification is asserted.
\end{abstract}
\noindent\textbf{Keywords:} Legendre family; character sums; Scholl representations; local monodromy; Artin conductors; parabolic cohomology.
\begin{center}\small
AI assistance and the scope of internal checking are disclosed in the revision note.\\
This preprint has not undergone independent specialist peer review.
\end{center}
\newpage
\setcounter{tocdepth}{1}
\begingroup\small\setlength{\parskip}{0pt}
\tableofcontents
\endgroup
\newpage
\section{Introduction and statements}
Character-twisted moments of the Legendre family give concrete finite sums
whose cohomological interpretation depends on boundary terms and arithmetic
normalizations. This paper treats cubic characters and the second and fourth
moments. Its central purpose is an explicit construction of the associated
representations and their Frobenius formulas, including the inert-prime sign. The consolidated local arguments now
also determine the two bad-prime factors and conductor exponents.
The exact finite-sum prescription underlying the numerical degree-four
recipe in the earlier preprint \cite{WalkeV1} is thereby identified
at good primes with a specified representation.

\subsection{Notation and Frobenius convention}
Put $K=\Q(\zeta_3)$, $E=\Q_2(\zeta_3)$ and
$U=\PP^1\setminus\{0,1,\infty\}$. Let
$\pi:\mathcal E\to U$ be the smooth projective Legendre family
\[
 \mathcal E_\lambda:y^2=t(t-1)(t-\lambda),\qquad
 \mathcal V=R^1\pi_*E.
\]
Let $\mathcal L_\rho$ be either nontrivial character summand of the actual
arithmetic cover $z^3=\lambda$. No additional constant Galois twist is
allowed in this definition. For an integer $r\in\{2,4\}$ set
\[
 H_r=H^1_!(U_{\overline K},\Sym^r\mathcal V\otimes\mathcal L_\rho),
 \qquad H^1_!=\operatorname{im}(H_c^1\longrightarrow H^1),
 \qquad W_2=H_2,\quad V_4=H_4.
\]
Frobenius always means \emph{geometric Frobenius on actual cohomology}.
Write $\tau$ for the character of $\Q_2(-1)$, so
$\tau(\Frob_{\mathfrak p})=N\mathfrak p$ and
$\tau=\chi_{\rm cyc}^{-1}$ in the usual convention on $\Q_2(1)$.
The modular comparison uses the cohomological dual of the conventional
arithmetic-Frobenius Deligne representation.

For a good residue field of cardinality $q$ and $\lambda\ne0,1$, put
\[
 a_q(\lambda)=q+1-\#\mathcal E_\lambda(\F_q),\qquad
 M_{3,j}(q)=\sum_{\lambda\in\F_q\setminus\{0,1\}}
                  \rho(\lambda)a_q(\lambda)^j.
\]
Here $\rho$ is the trace character of the specified Kummer summand.
When taking powers of one Frobenius, lift it by the finite-field norm.
At an inert rational prime of $K$, $\F_{p^2}$ is the initial residue
field; no nontrivial cubic character of $\F_p^\times$ is being lifted.
Define
\[
 v(q)=-M_{3,4}(q)+3qM_{3,2}(q)+q^2-1.
\]

\begin{theorem}[Summary of the fourth-moment results]\label{thm:main}
The representation $V_4$ has dimension four over $E$ and is unramified
outside the primes over $2,3$. At every good prime of $K$,
$\Tr(\Frob_{\mathfrak p}\mid V_4)=v(N\mathfrak p)$.
There are explicit four-dimensional $\Q_2$-representations $W_4^+$ and
$W_4^-$ of $G_\Q$, unramified outside $2,3$, whose restrictions to $G_K$
after scalar extension are $V_4$, with
\[
 W_4^-\simeq W_4^+\otimes\chi_{-3}.
\]
Both have a perfect alternating pairing with values in $\Q_2(-5)$
and are pure of weight five at every rational prime $p\ge5$.
Each is the $2$-adic realization of a rational Chow motive cut out on a
smooth projective fivefold, with Hodge type $(5,0)^2+(0,5)^2$;
see Theorem~\ref{thm:scholl-realization}.
For $p\ge5$, $p\equiv2\pmod3$, choose $d\in\F_p$ nonsquare,
$r^2=d$ in $\F_{p^2}$, and a primitive cubic character $\rho$ there.
Put
\[
 E'_x:Y^2=X(X^2-2xX+x^2-d),\quad a'_p(x)=p+1-\#E'_x(\F_p),
\]
\[
 S_p=\sum_{x\in\F_p}\rho(x+r)
       \bigl(a'_p(x)^4-3p\,a'_p(x)^2+p^2\bigr).
\]
Then $S_p\in\Z$ is independent of these choices and
\[
 \Tr(\Frob_p\mid W_4^+)=-S_p-1,\qquad
 \Tr(\Frob_p\mid W_4^-)=S_p+1.
\]
For $W_4^-$, set $c_1(p)=v(p)$ at split primes and $c_1(p)=S_p+1$
at inert primes, and set $c_2(p)=(c_1(p)^2-v(p^2))/2$ with the
character convention above. Then
\[
 \det(1-\Frob_pT\mid W_4^-)
 =1-c_1(p)T+c_2(p)T^2-p^5c_1(p)T^3+p^{10}T^4.
\]
\end{theorem}
\begin{theorem}[Local data and strong irreducibility]\label{thm:main-local}
Let $W_{4,\ell}^{\pm}$ be the realizations defined by the same rational
projectors as $W_4^{\pm}$ in Theorem~\ref{thm:scholl-realization}.
At $2$, for every $\ell\ne2$, both have zero inertia invariants and
Swan conductor zero. At $3$, for every $\ell\ne3$, the plus descent
has a one-dimensional invariant space with geometric Frobenius $3$,
whereas the minus descent has no invariants; both Swan conductors vanish.
Consequently the local Euler polynomials and Artin exponents are
\[
\begin{array}{c|cc|cc|c}
 &P_2(T)&a_2&P_3(T)&a_3&N\\\hline
W_4^+&1&4&1-3T&3&432\\
W_4^-&1&4&1&4&1296
\end{array}
\]
Here $N$ is the product of the local conductor powers, using a coefficient
prime different from the residue characteristic for each local calculation.
For every $\ell\ne2,3$, it is the prime-to-$\ell$ Artin conductor of
$W_{4,\ell}^{\pm}$. The displayed local data are independent of $\ell$
within their stated ranges.
For every $\ell\ne3$, both realizations are absolutely irreducible and
remain so after restriction to any finite number-field extension of $\Q$.
\end{theorem}
The local assertions are proved in Sections~\ref{sec:local-three} and
\ref{sec:local-two}, the conductor deduction in Section~\ref{sec:conductors},
and strong irreducibility in Section~\ref{sec:irreducibility}.
The last assertion includes the original $2$-adic realizations; no
irreducibility claim at coefficient prime $3$ is made here.

The construction, pairing and inert calculation are proved in
Sections~\ref{sec:geometry}, \ref{sec:descent}, \ref{sec:scholl} and \ref{sec:inert}.
The cubic second-moment comparison is treated separately in
Section~\ref{sec:cubic}; its modular conclusion uses the explicit exact
finite checks stated and independently recomputed there. These rank-two and rank-four objects
must not be conflated.

\subsection{Scope of this revision}
This consolidated research draft continues the version-2 preprint
\cite{WalkeV2}, with its earlier predecessor \cite{WalkeV1} retained for
provenance. It brings the local-at-$3$, dyadic tameness, dyadic invariant,
and strong-irreducibility arguments into the same manuscript.
The earlier blanket open obligations for these local data are superseded
by Theorem~\ref{thm:main-local}. Section~\ref{sec:claims} records the
current proof dependencies and unresolved questions. This draft has
internal AI-assisted review; external specialist review remains outstanding.
It does not certify the archived floating-point tables or promote their
functional-equation sign and special values to theorems.

The supplementary corrections retain known trace-formula consequences,
the Jacobi-sum normalization, and an analytic period substitution. Those
statements are explicitly distinguished from the main construction and
from numerical special-value observations. No classical elliptic BSD rank
or leading-term theorem is asserted.

\section{Relation to earlier work}\label{sec:related-work}

The cubic cover used here belongs to an established family.
Kazalicki--Scholl \cite[Sections~7--8]{KS} describe the genus-zero
groups $\Phi_0(N)$ and the elliptic family
$Y^2=X(X+1)(X+t^N)$. For $N=3$ and $t=z$, substituting $x=-X$
in our equation gives
\[
 y^2=-X(X+1)(X+z^3).
\]
Thus the two families differ by the constant quadratic twist $-1$.
The corresponding $\operatorname{Sym}^4$ coefficient systems are
isomorphic over $\mathbb Q$, since the fourth power of the twisting
character is trivial. Their analyzed weight-three representation uses
$R^1$; it does not by this comparison identify the fourth-symmetric-power constituent
studied here. Neither the underlying family nor its cubic-cover geometry
is claimed as new.

Scholl's construction places parabolic cohomology of symmetric powers
in the general theory of motives attached to noncongruence cusp forms;
see \cite[Introduction]{Scholl}. More precise projector references are
\cite[Sections~1.1--1.2]{SchollMotives} and
\cite[Sections~2.2--2.8 and 4.2--4.4]{SchollVanishing}.
Our local calculations use the standard
etale-cohomological framework. Section~\ref{sec:scholl} applies these
established constructions with an explicit full-level cover and projectors
matching the literal normalized deck and inversion operators. The degree-four
representations treated by Atkin--Li--Liu--Long \cite{ALLL} provide
further relevant context. Their automorphy theorem requires quaternion
multiplication satisfying specified compatibility conditions. We have
not established those conditions for $W_4^+$ or $W_4^-$
\cite[Definition~3.1.1 and Theorem~4.3.1]{ALLL}. In fact
$J=A-A^2$ exchanges the two spaces, while $I$ is scalar on each;
these operators are not a quaternion-multiplication pair on either
four-dimensional summand.

The auxiliary moment identities also have substantial precedents.
The quadratic-twist corollaries use the Legendre Hecke trace formulas
recorded in \cite[Theorem~2.3]{Grove}, with their attribution to Ahlgren
and Ahlgren--Ono \cite{Ahlgren,AhlgrenOno}. The weight-four forms used in the
second-moment comparisons occur in Long--Tu--Yui--Zudilin
\cite[Theorem~2 and Table~1]{LTYZ}; those modular forms and the
associated hypergeometric modularity are not discoveries of this note.

There is also an exact parameter overlap with Allen--Grove--Long--Tu
\cite{EHMM2}. The datum in our cubic second-moment normalization,
\[
 \left\{\begin{matrix}
 \tfrac12,\tfrac12,\tfrac13,\tfrac13\\
 1,1,\tfrac56,\tfrac56
 \end{matrix}\right\},
\]
is their $\mathrm{HD}_4(j/12)$ at $j=4$, defined in the published
article's introduction immediately before equation~(1.1).
(The corresponding definition is equation~(1.2) in arXiv v1.)
Their Theorem~1.4 concerns its specialization at argument $-1$;
our normalization identity uses argument $1$ and an explicit squared
Jacobi-sum factor. Equality of the parameter lists alone does not
identify these specializations or identify either with $W_4^\pm$.

The results established here specify the arithmetic coefficient sheaf,
construct its two rational descents, and derive their good-prime traces,
including the inert-prime boundary sign. These statements concern the
precisely normalized objects defined in the paper. Their exact priority
has not been established by the literature comparison, and no priority
claim is made for the constituent or trace formula. The conductor calculation in Sections~\ref{sec:local-three}--\ref{sec:conductors}
is an independent local argument, not a consequence of this comparison.
No automorphic identification or classical elliptic BSD consequence
follows from the literature overlap.

\section{Geometry and trace formulas over K}\label{sec:geometry}

Fix compatible primitive cube roots in
\(K=\Q(\zeta_3)\) and \(E=\Q_2(\zeta_3)\), and put
\(\mathcal O=\Z_2[\zeta_3]\),
\(S=\operatorname{Spec}\mathcal O_K[1/6]\), and
\(U=\PP^1_S\setminus\{0,1,\infty\}\).
Let \(\pi:\mathcal E\to U\) be the smooth projective Legendre elliptic
scheme, with its zero section and affine equation
\[
 \mathcal E_\lambda:\quad y^2=x(x-1)(x-\lambda),
 \qquad \mathcal V=R^1\pi_*E.
\]
The cubic torsor \(f:T\to U\), \(z^3=\lambda\), has deck
transformation \(a(z)=\zeta_3z\). We use pullback as the deck action.
Let \(\mathcal L_\rho\) be the \(\zeta_3\)-eigensheaf of \(a^*\)
in \(f_*E\). This definition fixes the arithmetic normalization;
in particular no constant Galois twist is inserted. For \(r\in\{2,4\}\)
set
\[
 \mathcal G_r=\Sym^r\mathcal V\otimes\mathcal L_\rho,
 \qquad
 H_r=H^1_!(U_{\bar K},\mathcal G_r)
 :=\operatorname{im}\bigl(H_c^1(U_{\bar K},\mathcal G_r)
                  \longrightarrow H^1(U_{\bar K},\mathcal G_r)\bigr).
\]
Thus \(W_2=H_2\) and \(V_4=H_4\). The subscript records the symmetric
power, and is distinct from the order three of the Kummer character.

All Frobenius operators below are geometric. Write \(\tau\) for the
character of \(E(-1)\), so \(\tau(\Frob_{\mathfrak p})=N\mathfrak p\).
For a prime \(\mathfrak p\nmid6\), put \(q=N\mathfrak p\) and let
\(\rho_q:\F_q^\times\to E^\times\) be the trace character of the
specified eigensheaf. For even \(k\) define
\[
 a_q(\lambda)=q+1-\#\mathcal E_\lambda(\F_q),\qquad
 M_{3,k}(q)=\sum_{\lambda\in\F_q\setminus\{0,1\}}
                  \rho_q(\lambda)a_q(\lambda)^k.
\]
The singular parameters are excluded throughout these moment sums.

\begin{theorem}\label{thm:common-geometry}
For \(r=2,4\), the representation \(H_r\) is continuous of dimension
\(r\) over \(E\), and is unramified outside the primes of \(K\) above
\(2,3\). There is an exact sequence of \(G_K=\Gal(\bar K/K)\)-representations
\begin{equation}\label{eq:common-full-sequence}
 0\longrightarrow E\longrightarrow H_c^1(U_{\bar K},\mathcal G_r)
 \longrightarrow H^1(U_{\bar K},\mathcal G_r)
 \longrightarrow E(-(r+1))\longrightarrow0.
\end{equation}
Both middle spaces have dimension \(r+1\), and their common parabolic
image is \(H_r\). At every \(\mathfrak p\nmid6\),
\begin{align}
 \Tr(\Frob_{\mathfrak p}\mid W_2)
   &=-M_{3,2}(q)-q-1,\label{eq:trace-W2}\\
 \Tr(\Frob_{\mathfrak p}\mid V_4)
   &=-M_{3,4}(q)+3qM_{3,2}(q)+q^2-1.
   \label{eq:trace-V4}
\end{align}
These assertions include the degree-two prime ideals of \(K\).
\end{theorem}

\subsection{Local monodromy and the arithmetic boundary}

At \(\lambda=0,1\), the discriminant
\(16\lambda^2(1-\lambda)^2\) has valuation two, while
\(c_4=16(\lambda^2-\lambda+1)\) is a unit. Thus the fibres are
multiplicative, with nontrivial unipotent inertia on \(\mathcal V\).
At infinity, the substitutions \(t=1/\lambda\), \(x=\lambda X\),
\(y=\lambda^2Y\) give
\[
 Y^2=tX(X-1)(X-t).
\]
This is the ramified quadratic twist of the multiplicative family at
\(t=0\). The geometric inertia types on \(\mathcal V\) at
\(0,1,\infty\) are consequently \(J_2(1),J_2(1),J_2(-1)\);
see also \cite[Chapter~15]{Katz}.

For clarity, the symmetric-power calculation is made over the
characteristic-zero field \(E\). In a local basis choose
\(N(e)=0\), \(N(h)=ce\), \(c\ne0\), for the unipotent logarithm.
The induced logarithm satisfies
\[
 N_r(e^{r-j}h^j)=jc\,e^{r-j+1}h^{j-1},\qquad
 N_r^r(h^r)=r!c^re^r\ne0.
\]
It therefore has one block of length \(r+1\). Since \(r\) is even,
the scalar \(-1\) at infinity disappears. No assertion that \(r!\)
is an integral unit is needed. With compatible tame generators the
Kummer scalars at zero and infinity are a primitive cube root \(\eta\)
and \(\eta^{-1}\), respectively. Hence
\begin{center}
\begin{tabular}{c|c|c|c}
 point & \(\mathcal V\) & \(\mathcal G_r\) &
          \(\dim_E\mathcal G_r^{I}\)\\ \hline
 \(0\) & \(J_2(1)\) & \(\eta J_{r+1}(1)\) & \(0\)\\
 \(1\) & \(J_2(1)\) & \(J_{r+1}(1)\) & \(1\)\\
 \(\infty\) & \(J_2(-1)\) & \(\eta^{-1}J_{r+1}(1)\) & \(0\)
\end{tabular}
\end{center}
The same calculation holds in every good geometric fibre. The integral
tameness argument below verifies that no Swan terms occur there.

\begin{lemma}\label{lem:arithmetic-boundary}
The inertia invariant and coinvariant lines at \(\lambda=1\) are
arithmetically
\[
 \mathcal G_r^{I_1}\simeq E,
 \qquad (\mathcal G_r)_{I_1}\simeq E(-r).
\]
In particular geometric Frobenius acts on the invariant line by \(1\),
both over \(K\) and at every good residue field.
\end{lemma}

\begin{proof}
Put \(v=x-1\) in the singular fibre at \(\lambda=1\). Its equation
\(y^2=(1+v)v^2\) has the two rational tangents \(y=\pm v\), so its
multiplicative reduction is split. The Tate torsion sequence, followed by
inverse limit and duality, gives the cohomological sequence
\[
 0\longrightarrow E\longrightarrow\mathcal V
      \longrightarrow E(-1)\longrightarrow0
\]
on the local punctured trait; see \cite[pp.~18--19]{Tate}.
Its nontrivial unipotent inertia fixes the first line. In its \(r\)-th
symmetric power, the invariant line is the \(r\)-th power of this
trivial line, whereas the coinvariant quotient is the \(r\)-th power
of the quotient \(E(-1)\). Finally the torsor fibre over \(1\) is
\(\mu_3\). It is rational over \(K\) and over each residue field at
\(\mathfrak p\nmid6\), so its chosen character line is arithmetically
trivial. Tensoring by that line gives both assertions.
\end{proof}

\subsection{Cohomology and the two boundary terms}

In this subsection cohomology is first taken over \(\bar K\); the same
arguments apply over good algebraically closed residue fields.
The absence of \(I_0\)-invariants implies
\(H^0(U,\mathcal G_r)=0\). It also implies
\(H^0(U,\mathcal G_r^\vee(1))=0\), since a constant Tate twist is
trivial on geometric inertia. A lisse section on a connected nonproper
smooth curve cannot have proper support unless it is zero. Poincare
duality \cite[Theorem~24.1]{Milne} therefore gives
\begin{equation}\label{eq:common-vanishing}
 H_c^0(U,\mathcal G_r)=H_c^2(U,\mathcal G_r)=0.
\end{equation}
Comparison with the complex thrice-punctured sphere, which retracts
onto a graph with one vertex and two loops, gives ordinary Euler
characteristic \(-(r+1)\) for a rank-\((r+1)\) local system.
Duality gives the same compact-support Euler characteristic. The
specialization below transports the dimensions to the good fibres.
In particular
\(\dim_E H_c^1(U,\mathcal G_r)=r+1\).

Let \(j:U_{\bar K}\hookrightarrow\PP^1_{\bar K}\), and let
\(i_1\) be the inclusion of the point \(1\). Lemma~\ref{lem:arithmetic-boundary}
and the local invariant table give an exact sequence of ordinary etale
sheaves
\[
 0\longrightarrow j_!\mathcal G_r\longrightarrow j_*\mathcal G_r
     \longrightarrow i_{1*}E\longrightarrow0.
\]
Its cohomology sequence and the vanishing of \(H^0(U,\mathcal G_r)\)
give
\begin{equation}\label{eq:common-left-boundary}
 0\longrightarrow E\longrightarrow H_c^1(U,\mathcal G_r)
   \longrightarrow H^1(\PP^1,j_*\mathcal G_r)\longrightarrow0.
\end{equation}
The edge map for Leray injects
\(H^1(\PP^1,j_*\mathcal G_r)\) into \(H^1(U,\mathcal G_r)\).
The composite of this injection with
\eqref{eq:common-left-boundary} is the natural compact-support map.
It follows that
\[
 H^1(\PP^1,j_*\mathcal G_r)=H_r,
 \qquad \dim_EH_r=r.
\]
These identifications and exact sequences are Galois equivariant.

To determine the right boundary, the same sheaf sequence and
\eqref{eq:common-vanishing} imply
\(H^2(\PP^1,j_*\mathcal G_r)=0\). Leray now gives a surjection from
\(H^1(U,\mathcal G_r)\) onto
\(\bigoplus_{s=0,1,\infty}H^1(I_s,\mathcal G_r)\), with kernel
\(H_r\). At \(0,\infty\), the operator \(T_s-1\) is invertible,
so the corresponding local groups vanish. At \(1\), tame inertia
cohomology identifies the local group with
\[
 H^1(I_1,\mathcal G_r)
   \simeq (\mathcal G_r)_{I_1}(-1)
   \simeq E(-(r+1)).
\]
The additional twist comes from the dual of the tame factor
\(\Z_2(1)\); geometric Frobenius acts on this factor of the local
cohomology by \(q\). Combining these observations proves
\eqref{eq:common-full-sequence}, including its final surjectivity, and
\(\dim_EH^1(U,\mathcal G_r)=r+1\).

\subsection{Integral tame specialization}

We give the specialization argument explicitly because \(U\) is not
proper. The eigensheaf has an integral lattice
\(\mathcal L_{\rho,\mathcal O}\): the order-three projector is
integral since \(3\) is invertible in \(\mathcal O\). Put
\[
 \mathcal G_{r,\mathcal O}
  =\bigl(\Sym^r(R^1\pi_*\Z_2)\otimes_{\Z_2}\mathcal O\bigr)
       \otimes_{\mathcal O}\mathcal L_{\rho,\mathcal O},
 \qquad \mathcal G_{r,m}=\mathcal G_{r,\mathcal O}/2^m.
\]
The finite monodromy of \(\mathcal G_{r,m}\) factors through the image
of a subgroup of
\(\operatorname{GL}_2(\Z/2^m)\times\mu_3\), whose order is
\[
 18\cdot2^{4(m-1)}.
\]
Thus for every residue characteristic \(p\ge5\), wild pro-\(p\)
inertia at a boundary point acts trivially. In characteristic zero there
is no wild inertia. This verifies relative tameness in the fibrewise
discrete-valuation sense of \cite[Expose~XIII, Definition~2.1.1]{SGA1}.
The sheaves are already lisse on \(U\).

Apply \cite[Expose~XIII, Corollary~2.9]{SGA1} to
\(\PP^1_S\), its relative normal-crossing boundary consisting of the
three disjoint sections, and the additive coefficient sheaf
\(\mathcal G_{r,m}\). The compactification is smooth, proper, and of
finite presentation. The coefficient sheaf is finite, locally constant,
constructible, relatively tame, and \(2\)-primary, with \(2\)
invertible on \(S\). The theorem supplies local constancy and
base change in degrees at most one. The same argument applies to the
integral dual coefficient systems and their Tate twists; their additional
finite cyclotomic factors have \(2\)-power order.

On a geometric fibre, the finite \(H^0\) groups form a
Mittag--Leffler system. The continuous-cohomology comparison
\cite[Sections~18--21]{Milne} consequently gives
\[
 H^1(U,\mathcal G_r)
   =\left(\varprojlim_m H^1(U,\mathcal G_{r,m})\right)
       \otimes_{\mathcal O}E.
\]
The finite-level specialization isomorphisms therefore pass to the
\(E\)-adic groups. This argument does not require finite-level
cohomology to be free.

For compact supports we use relative duality, rather than applying
smooth proper base change to the open curve. The morphism \(g:U\to S\)
is smooth, separated, compactifiable, and of relative dimension one;
\(S\) is quasi-compact and quasi-separated, and every \(2^m\) is
invertible on it. The relative trace and smooth-duality theorems
\cite[Expose~XVIII, Theorem~2.9, Var.~2, and Theorem~3.2.5]{SGA4}
apply. Their base-change compatibility, followed by passage to \(E\),
identifies
\[
 H_c^1(U,\mathcal G_r)
       \simeq H^1(U,\mathcal G_r^\vee(1))^\vee
\]
compatibly with specialization. Moreover the natural transformation
\(Rg_!\to Rg_*\) is compatible with these base-change maps.
Accordingly the specialization square with horizontal maps
\(H_c^1\to H^1\) commutes, and both vertical maps are isomorphisms.
Taking images over \(E\) gives an isomorphism of the parabolic
cohomology groups, equivariant for residue-field Galois action.

The compatible finite systems give a continuous representation factoring
through \(\pi_1(S)\). Their parabolic image is a stable subquotient,
so \(H_r\) is unramified outside \(2,3\). At a good prime its
Frobenius is identified with the Frobenius of the special-fibre
parabolic group. This proves precisely the specialization needed for
the arithmetic trace calculation.

\subsection{The two trace formulas}

Let \(\alpha,\beta\) be the geometric Frobenius eigenvalues of a
smooth Legendre fibre over \(\F_q\). Thus
\(\alpha+\beta=a_q(\lambda)\) and \(\alpha\beta=q\). The
symmetric-power recurrence gives
\[
 \Tr(\Frob_\lambda\mid\Sym^2\mathcal V)
      =a_q(\lambda)^2-q,
 \qquad
 \Tr(\Frob_\lambda\mid\Sym^4\mathcal V)
      =a_q(\lambda)^4-3qa_q(\lambda)^2+q^2.
\]
The compact-support trace formula \cite{StacksTrace}, the vanishings
\eqref{eq:common-vanishing}, and the arithmetic boundary sequence
\eqref{eq:common-left-boundary} therefore give
\[
 \Tr(\Frob_q\mid H_r)
   =-\sum_{\lambda\ne0,1}
          \Tr(\Frob_\lambda\mid\mathcal G_{r,\lambda})-1.
\]
Since \(\sum_{\lambda\ne0,1}\rho_q(\lambda)=-1\), this is
\eqref{eq:trace-W2} for \(r=2\), and \eqref{eq:trace-V4} for
\(r=4\). The final \(-1\) is the trace of the actual boundary
line, not a correction inferred from its dimension. The preceding
specialization identifies these special-fibre traces with the
\(G_K\)-representations in Theorem~\ref{thm:common-geometry}, completing
its proof.

For powers of a fixed Frobenius, the character on \(\F_{q^d}^\times\)
must be \(\rho_q\circ N_{\F_{q^d}/\F_q}\). With this convention the
same formulas compute \(\Tr(\Frob_{\mathfrak p}^{d}\mid H_r)\).
At an inert rational prime \(p\equiv2\pmod3\), the initial residue
field is \(\F_{p^2}\); its cubic character is not a nontrivial norm
lift from \(\F_p\). These formulas determine the trace over \(K\)
and do not yet choose a trace at \(p\) for an extension to \(G_\Q\).

Parameter inversion gives
\(a_q(1/\lambda)=\chi_q(\lambda)a_q(\lambda)\), because the two
elliptic curves differ by the indicated quadratic twist. For even
\(k\), changing variables in the moment consequently yields
\[
 M_{3,k}(q;\rho_q)=M_{3,k}(q;\rho_q^{-1})\in\Z.
\]
Indeed the sum is a cyclotomic integer fixed by conjugation. This
character independence does not itself supply coefficient descent.

Section~\ref{sec:descent} places both cubic eigenspaces in an ambient
\(\Q_2\)-representation \(H_{\mathrm{nt}}\), constructs the
arithmetic inversion operator \(I\), and uses it to define the two
extensions \(W_4^\pm\) to \(G_\Q\).

\section{The cubic second moment and its modular comparison}\label{sec:cubic}
Throughout this section the character has order three. Put
\[
 \mathcal G_2=\Sym^2\mathcal V\otimes\mathcal L_\rho,
 \qquad W_2=H^1_!(U_{\overline K},\mathcal G_2).
\]
The preceding geometric construction gives $\dim_EW_2=2$, unramifiedness
outside the primes above $2,3$, and, for every other prime ideal
$\mathfrak p$ with $q=N\mathfrak p$,
\begin{equation}\label{eq:cubic-trace}
 \Tr(\Frob_{\mathfrak p}\mid W_2)=-M_{3,2}(q)-q-1.
\end{equation}
Here Frobenius is geometric, $\tau=E(-1)$ has value $q$, and the Kummer
normalization is the actual torsor $z^3=\lambda$. We first obtain the
remaining determinant constraint geometrically. The modular comparison
then uses the separately recorded finite arithmetic certificate.

\subsection{Arithmetic inversion and the determinant}
Let $i(\lambda)=\lambda^{-1}$. On the etale cover $u^2=\lambda$, the map
\[
 \phi_u:\mathcal E_\lambda\longrightarrow\mathcal E_{1/\lambda},
 \qquad (x,y)\longmapsto (x/\lambda,y/u^3)
\]
is an elliptic-curve isomorphism: substitution gives
$(y/u^3)^2=(x/\lambda)(x/\lambda-1)(x/\lambda-1/\lambda)$.
Replacing $u$ by $-u$ composes with $[-1]$, whose action on
$\Sym^2\mathcal V$ is trivial. Thus the induced isomorphism
$i^*\Sym^2\mathcal V\simeq\Sym^2\mathcal V$ descends arithmetically.
The map $z\mapsto z^{-1}$ on the cubic cover inverts the deck group,
so $i^*\mathcal L_\rho\simeq\mathcal L_\rho^\vee$, $E$-linearly over $K$.
This step does not conjugate the coefficient field.

The elliptic cup product induces the nondegenerate symmetric pairing
\[
 (v_1v_2,w_1w_2)\longmapsto
 \langle v_1,w_1\rangle\langle v_2,w_2\rangle+
 \langle v_1,w_2\rangle\langle v_2,w_1\rangle
 \quad\text{with values in }E(-2).
\]
Consequently $i^*\mathcal G_2\simeq\mathcal G_2^\vee(-2)$.
Poincare duality \cite[Theorem 24.1]{Milne} induces a perfect pairing
on parabolic images: the kernels of the maps $H_c^1\to H^1$ are the
annihilators of the corresponding dual images. Hence
\[
 H^1_!(U,\mathcal G_2^\vee)\simeq W_2^\vee(-1),
 \qquad W_2\simeq W_2^\vee(-3).
\]
The inversion pullback commutes with $G_K$ and with the compact-support
map, so these are arithmetic isomorphisms. Taking determinants gives
\begin{equation}\label{eq:cubic-det-sign}
 \det W_2=\tau^3\varepsilon,\qquad \varepsilon^2=1,
\end{equation}
where $\varepsilon$ is continuous and unramified outside $2,3$.
In usual cyclotomic notation this is
$\chi_{\rm cyc}^{-3}\varepsilon$. The assertion
$\varepsilon=1$ will use the finite comparison below.

\subsection{The modular representation and residual conditions}
Let $f=f_{36}$ be the weight-four normalized newform of level $36$
and trivial nebentypus, labelled $36.4.a.a$. Its eta expression is
\cite[Table 1]{LTYZ}
\begin{equation}\label{eq:f36-eta}
 f(z)=\frac{\eta(6z)^{14}}{\eta(2z)^3\eta(18z)^3}
       -3\eta(2z)^3\eta(6z)^2\eta(18z)^3.
\end{equation}
Use the cohomological representation $R_f$ over $E$: it is the dual of
the conventional Deligne representation normalized at arithmetic
Frobenius; see the modular-form motive and its etale realization in
\cite[Theorem~1.2.4(i)]{SchollMotives}.
Thus geometric Frobenius at $p\ge5$ on $R_f$ has trace
$a_f(p)$ and determinant $p^3$. Both $W_2$ and $R_f|_{G_K}$ are
unramified outside primes above $2,3$.

Their compact images admit stable $\mathcal O$-lattices, with
$\mathcal O/(2)=\F_4$. Full rational $2$-torsion on every smooth
Legendre fibre over $\F_q$ makes $a_q(\lambda)$ even. Therefore
$M_{3,2}(q)\in4\mathcal O$, and
\eqref{eq:cubic-trace}--\eqref{eq:cubic-det-sign} give
\[
 \Tr W_2\equiv0\pmod2,\qquad \det W_2\equiv1\pmod2.
\]
Initially these congruences hold on good Frobenius classes;
Chebotarev density and continuity extend them to $G_K$.

The carried certificate records, and the independent calculation described
below rechecks,
$f\equiv\eta(3z)^8\pmod2$ through the Sturm bound
\[
 \frac4{12}[\operatorname{SL}_2(\Z):\Gamma_0(36)]
 =\frac4{12}\,72=24.
\]
Here $\eta(3z)^8\in S_4(\Gamma_0(9))\subset S_4(\Gamma_0(36))$
by \cite[Table~1]{LTYZ}. Both sides have integral Fourier coefficients.
Sturm's theorem \cite[Theorem~1, p.~276]{Sturm} makes this a full
Fourier-series congruence,
subject to that recorded finite check. With $Q_z=e^{2\pi iz}$,
\[
 \eta(3z)^8\equiv\eta(24z)
 =\sum_{m\ge1}\chi_{12}(m)Q_z^{m^2}\pmod2.
\]
The congruence is the characteristic-two product identity; the equality
is Euler's pentagonal-number identity, with
$\chi_{12}(m)=(12/m)$ the Kronecker symbol. Prime-index coefficients of $f$ are therefore
even. At a split prime of $K$, the restricted modular trace is $a_f(p)$;
at an inert prime it is $a_f(p)^2-2p^3$, again even. Its determinant
is $q^3$. Density gives the required global congruences for
$R_f|_{G_K}$ as well. The residue field $\F_4$ is allowed in the
comparison criterion used below.

\subsection{The carried finite certificate and its consequence}
Let $K_S$ be the compositum of the quadratic extensions of $K$
unramified outside the primes above $2,3$. The Eisenstein integer ring
has class number one. Removing square factors from an allowed
squareclass therefore leaves an $S$-unit. Its unit squareclass is
generated by $-1$, and the primes above $2$ and $3$ can be generated
by $2$ and $\sqrt{-3}$. Their valuations, together with the nonsquare
unit $-1$, prove independence. Consequently
\[
 K_S=K\bigl(\sqrt{-1},\sqrt2,\sqrt{\sqrt{-3}}\bigr),
 \qquad\Gal(K_S/K)\simeq\F_2^3.
\]
The seven nonzero vectors constitute a non-cubic set: a homogeneous
cubic polynomial function vanishing on them also vanishes at zero,
and hence vanishes on the whole vector space.

Table~\ref{tab:cubic-certificate} reproduces the cubic rows of the
carried exact certificate. A prime ideal is specified by $p$ and the
image $r$ of $\sqrt{-3}$ modulo it. A bit is $0$ for a square and $1$
for a nonsquare; the displayed order is
$(\sqrt{-3},2,-1)$, the reverse of the basis above.
\begin{table}[htbp]
\centering
\begin{tabular}{rrrrr}
\hline
$p$ & $r$ & bits & $t=a_f(p)$ & $d$ \\
\hline
7  & 2  & 001 & 8    & 343 \\
7  & 5  & 101 & 8    & 343 \\
13 & 6  & 110 & $-10$  & 2197 \\
19 & 4  & 011 & $-100$ & 6859 \\
19 & 15 & 111 & $-100$ & 6859 \\
37 & 16 & 010 & $-226$ & 50653 \\
73 & 17 & 100 & 26   & 389017 \\
\hline
\end{tabular}
\caption{Previously recorded exact cubic comparisons, with $d=p^3$.
The seven rows and Sturm comparison have also been recomputed by a new
standalone exact implementation.}
\label{tab:cubic-certificate}
\end{table}

The certificate consists of the supplied
\texttt{V2\_01\_exact\_livne.py}, its exact arithmetic helpers,
\texttt{V2\_exact\_livne.txt}, and
\texttt{exact\_livne\_certificate.json}. Static inspection shows direct
finite-field character sums and integer eta-series arithmetic for the
cubic comparisons, with no FFT rounding. The certificate uses the
norm-lifted character over $\F_{p^2}$. Thus, by the geometric trace
formula,
\begin{align*}
 t&=-M_{3,2}(p)-p-1,\\
 u&=-M_{3,2}(p^2)-p^2-1=\Tr(\Frob_{\mathfrak p}^{2}\mid W_2),\\
 d&=\frac{t^2-u}{2}.
\end{align*}
The inversion identity for even Legendre moments makes the two
primitive cubic-character choices equivalent here. The recorded
values satisfy $t=a_f(p)$ and $d=p^3$ at every listed prime; the
recorded Sturm comparison reaches degree $24$.
\emph{The archived programs were not rerun.} A new standalone implementation,
\texttt{independent\_cubic\_exact.py}, independently recomputes the
displayed cubic comparisons using direct finite-field character sums and
integer eta-series arithmetic. It uses only the Python standard library;
the mathematical calculations use arbitrary-precision integers, without
FFT or floating-point arithmetic. Over $\F_{p^2}=\F_p[s]/(s^2-d)$ it uses
$\chi_{p^2}(a+bs)=\chi_p(a^2-db^2)$ and the norm-lifted cubic character.
The new log records agreement at the five rational primes underlying the
seven rows, all seven square-class vectors, and the congruence through
degree $24$. Historical expected values are read only after the new
calculation, for comparison. Checks include every rational-fibre quadratic
base-change identity and a second direct multiplication calculation over
$\F_{49}$. Source hashes and the new log are supplied in the release
supplement. This rechecks the bounded arithmetic; it does not re-prove the
geometric hypotheses, newform identification, or external comparison
theorems. The theorem below retains its explicit finite-check dependencies.

\begin{theorem}[Cubic modular comparison using the carried certificate]
\label{thm:cubic-modular}
Assuming the finite exact comparisons and the Sturm comparison just
specified, the representations have isomorphic semisimplifications:
\[
 W_2^{\rm ss}\simeq(R_f|_{G_K})^{\rm ss},
 \qquad \det W_2=\tau^3.
\]
For every good prime ideal $\mathfrak p$ of $K$,
\[
 M_{3,2}(N\mathfrak p)
 =-\Tr(\Frob_{\mathfrak p}\mid R_f|_{G_K})-N\mathfrak p-1.
\]
In particular, for $p\ge5$ and $p\equiv1\pmod3$,
$M_{3,2}(p)=-a_f(p)-p-1$; for $p\equiv2\pmod3$,
\[
 M_{3,2}(p^2)=-a_f(p)^2+2p^3-p^2-1.
\]
\end{theorem}
\begin{proof}
Livne's criterion in \cite[Theorem 4.1]{DGP} applies to continuous
rank-two representations over a finite extension of $\Q_2$.
The common ramification set, vanishing residual traces, and equal
residual determinants were established above. The seven rows give
exact trace and determinant agreement on a non-cubic set in
$\Gal(K_S/K)$. The criterion gives the asserted semisimplification.
If it is formulated using the opposite Frobenius convention, invert
both matrices: $\Tr(A^{-1})=\Tr(A)/\det A$, and the quadratic test
classes are unchanged under inversion.

Alternatively to taking the determinant of the comparison,
\eqref{eq:cubic-det-sign} and the table give
$\varepsilon(\Frob_{\mathfrak p})=d/p^3=1$ on all seven nonzero
classes. Since $\varepsilon$ factors through $\Gal(K_S/K)$, it is
trivial. Finally apply \eqref{eq:cubic-trace}; restriction at an inert
rational prime replaces Frobenius by its square.
\end{proof}

This section proves the stated cubic second-moment comparison with its
explicit certificate dependency. It supplies no geometric proof for
characters of orders four or six, and no bad-prime Euler-factor or
rank-four modularity assertion.

\section{Arithmetic involution, descent and polarization}\label{sec:descent}

We construct two four-dimensional representations over $\Q_2$ whose
restrictions to $G_K$, after extending coefficients to $E$, are $V_4$.
Throughout this section $\chi=\chi_{-3}$ is the quadratic character of
$K/\Q$, and $\tau$ denotes the character of $\Q_2(-1)$.

\subsection{The rational ambient representation}

Over $\Q$, let
\[
 f:T=\PP^1_z\setminus\{0,\infty,\mu_3\}\longrightarrow U,
 \qquad \lambda=z^3,
 \qquad \mathcal A_0=\Sym^4(R^1\pi_*\Q_2).
\]
The divisor $\mu_3$ and the finite etale map $f$ are defined over $\Q$.
Set
\[
 H=H^1_!(T_{\overline\Q},f^*\mathcal A_0),\qquad
 A=a^*,\quad a(z)=\zeta_Kz,\qquad
 e_{\mathrm{nt}}=\frac{2-A-A^2}{3},\quad
 H_{\mathrm{nt}}=e_{\mathrm{nt}}H.
\]
Here $\zeta_K\in K$ is a chosen primitive cube root, and $A$ means
pullback, with the identity identification $a^*f^*\mathcal A_0=f^*\mathcal A_0$.
The operator $A$ is $\Q_2$-linear, although $a$ is defined over $K$.
If $R$ denotes the Galois action on $H$, conjugation of the deck map gives
\begin{equation}\label{eq:descent-deck-galois}
 A^3=1,\qquad R(g)AR(g)^{-1}=A^{\chi(g)}.
\end{equation}
Consequently $e_{\mathrm{nt}}$ and its image are $G_\Q$-stable.

Finite pushforward is exact, and the projection formula identifies
$f_*f^*\mathcal A_0$ with $\mathcal A_0\otimes f_*\Q_2$.
These assertions can be checked on the finite sums of geometric stalks;
see also \cite[Proposition 59.55.2, tag 03QP]{StacksFinite}.
Since $f$ is proper, the same identification applies to compactly
supported cohomology and commutes with $H_c^1\to H^1$.
Thus it also identifies their parabolic images. With the chosen
coefficient root $\zeta_E\in E$, take $\mathcal L_\rho$ to be the
$A=\zeta_E$ eigensheaf of $f_*E$, as in the preceding construction.
It follows that
\begin{equation}\label{eq:descent-eigenspaces}
 H_{\mathrm{nt}}\otimes_{\Q_2}E
   =V_{4,\rho}\oplus V_{4,\rho^{-1}},\qquad
 \dim_{\Q_2}H_{\mathrm{nt}}=8.
\end{equation}
The dimension follows from the previously proved dimension four of
each character summand. Both summands are unramified at primes of $K$
above $p\ge5$. Since $K/\Q$ is unramified at such $p$, the inertia
group of $\Q$ lies in $G_K$, so $H_{\mathrm{nt}}$ is also unramified there.

\subsection{An actual coefficient involution}

\begin{lemma}\label{lem:descent-involution}
There is a $\Q_2$-linear operator $I$ on $H_{\mathrm{nt}}$ such that
\[
 I^2=1,\qquad IR(g)=R(g)I,\qquad IAI^{-1}=A^{-1}.
\]
It is induced by an arithmetic coefficient isomorphism over inversion
$z\mapsto z^{-1}$ and respects the compact-support and boundary maps.
\end{lemma}

\begin{proof}
Put $i(\lambda)=1/\lambda$. On the finite etale double cover
$u^2=\lambda$ of $U$, the zero-preserving elliptic isomorphism
\begin{equation}\label{eq:descent-elliptic-map}
 \phi_u:\mathcal E_\lambda\longrightarrow\mathcal E_{1/\lambda},
 \qquad (x,y)\longmapsto(x/\lambda,y/u^3)
\end{equation}
is verified by
\[
 (x/\lambda)(x/\lambda-1)(x/\lambda-1/\lambda)
   =x(x-1)(x-\lambda)/\lambda^3=(y/u^3)^2.
\]
Replacing $u$ by $-u$ composes this map with $[-1]$, which acts as
$+1$ on $\Sym^4R^1\pi_*\Q_2$. The two pullbacks therefore agree
on the double-cover relation and descend to an isomorphism
\[
 \theta:i^*\mathcal A_0\longrightarrow\mathcal A_0
\]
over $\Q$; in fact all these maps are defined over $\Z[1/6]$.
For the second inversion choose $1/u$. Directly on both coordinates,
$\phi_{1/u}\circ\phi_u=\operatorname{id}$, whence
\begin{equation}\label{eq:descent-coefficient-cocycle}
 \theta\circ i^*\theta=\operatorname{id}.
\end{equation}
Choosing $-1/u$ gives the same identity on fourth symmetric powers.

Let $b(z)=1/z$. Since $fb=if$, pullback of $\theta$ gives
$\Theta:b^*f^*\mathcal A_0\to f^*\mathcal A_0$.
Define $I=\Theta_*b^*$ on cohomology. The map $b$ extends to an
automorphism of $\PP^1_z$ preserving its boundary as a set. Its
pullback and the coefficient isomorphism therefore act on $H_c^1$,
$H^1$, and the natural map between them; they act on the parabolic
image and on the boundary exact sequence as well. Equation
\eqref{eq:descent-coefficient-cocycle} gives $I^2=1$.
Both $b$ and $\Theta$ are defined over $\Q$, so $I$ commutes with $R$.
Finally $bab=a^{-1}$, while the coefficient map is pulled back from
the $\lambda$-line and the deck coefficient maps are identities.
This proves $IAI^{-1}=A^{-1}$ on coefficient systems and hence on
cohomology. It also proves that $I$ preserves $H_{\mathrm{nt}}$ and
exchanges the summands in \eqref{eq:descent-eigenspaces}.
\end{proof}

\subsection{The two descents and the quadratic twist}

\begin{theorem}\label{thm:descent-rational}
The spaces
\[
 W_4^+=H_{\mathrm{nt}}^{I=1},\qquad
 W_4^-=H_{\mathrm{nt}}^{I=-1}
\]
are continuous four-dimensional $\Q_2$-representations of $G_\Q$,
unramified outside $2,3$, and
\[
 (W_4^\pm|_{G_K})\otimes_{\Q_2}E\simeq V_4,
 \qquad W_4^-\simeq W_4^+\otimes\chi.
\]
\end{theorem}

\begin{proof}
First consider extension of the base field with coefficients $E$.
Let $V=\ker(A-\zeta_E)$ in $H_{\mathrm{nt}}\otimes E$, and put
$\delta(g)=(1-\chi(g))/2$. For $g\notin G_K$, both $R(g)$ and $I$
exchange the two eigenspaces. Hence
\begin{equation}\label{eq:descent-extensions}
 R_+(g)=I^{\delta(g)}R(g)|_V,\qquad
 R_-(g)=(-I)^{\delta(g)}R(g)|_V
\end{equation}
preserve $V$. The commutation in Lemma~\ref{lem:descent-involution}
and additivity of $\delta$ modulo two show that these maps are
representations. Their restriction to $G_K$ is the original action.

Coefficient descent is supplied by the displayed rational eigenspaces.
Indeed
\[
 \Phi_\pm:V\longrightarrow W_4^\pm\otimes E,
 \qquad v\longmapsto v\pm Iv
\]
are isomorphisms whose inverses are projection to $V$. For $g\notin G_K$,
\[
 R(g)(v\pm Iv)=\Phi_\pm(\pm IR(g)v);
\]
for $g\in G_K$ the intertwining is immediate. Thus these rational
spaces realize exactly \eqref{eq:descent-extensions} and have dimension
four. Continuity and the ramification assertion follow from $H_{\mathrm{nt}}$.

For an explicit coefficient-descent datum, let
$C=\operatorname{id}\otimes\sigma$, where
$\sigma\in\Gal(E/\Q_2)$ sends $\zeta_E$ to $\zeta_E^{-1}$.
This semilinear operator commutes with $A,I,R$ and exchanges the
two eigenspaces. On $V$ the operators $D_+=IC$, $D_-=-IC$ are
semilinear involutions commuting with the corresponding $R_\pm$.
Under $\Phi_\pm$ they are ordinary coefficient conjugation. This
checks the coefficient cocycle separately from base-field extension.

Finally put $J=A-A^2$ on $H_{\mathrm{nt}}$. Since $A+A^2=-1$ there,
\[
 J^2=-3,\qquad IJ=-JI,\qquad R(g)JR(g)^{-1}=\chi(g)J.
\]
Thus $J$ is an invertible $\Q_2$-linear map from $W_4^+$ to $W_4^-$,
intertwining the latter with $W_4^+\otimes\chi$.
\end{proof}

\begin{remark}
The signs refer to the literal map \eqref{eq:descent-elliptic-map}:
$W_4^+$ is its $+1$ eigenspace. Replacing $I$ by $-I$ gives $W_4^-$.
This constructs two particular extensions without any irreducibility
assumption or classification of all extensions.
\end{remark}

\subsection{Boundary action}

The boundary lines above $\lambda=1$ have trivial arithmetic coefficient
action. Formula \eqref{eq:descent-elliptic-map} extends regularly there,
since $\lambda$ and $u$ are units; at $u=1$ it specializes to the
identity. The choice $u=-1$ has the same effect on $\Sym^4$.
Consequently $I$ acts on the three boundary lines by the permutation
$z\mapsto z^{-1}$, with coefficient $+1$.
The nontrivial deck boundary is therefore the augmentation subspace
$B_{\mathrm{nt}}\subset\Q_2[\mu_3]$. The boundary contributions at
$0,\infty$ are killed by $e_{\mathrm{nt}}$, equivalently because
the two twisted eigensheaves have no inertia invariants there.
Writing $C_{\mathrm{nt}}=e_{\mathrm{nt}}H_c^1(T_{\overline\Q},f^*\mathcal A_0)$,
we obtain
\begin{equation}\label{eq:descent-boundary-exact}
 0\longrightarrow B_{\mathrm{nt}}\longrightarrow C_{\mathrm{nt}}
   \longrightarrow H_{\mathrm{nt}}\longrightarrow0.
\end{equation}
In $B_{\mathrm{nt}}$, the lines generated by
$2[1]-[\zeta_K]-[\zeta_K^2]$ and $[\zeta_K]-[\zeta_K^2]$
have $I$ eigenvalues $+1,-1$ and Galois characters $1,\chi$, respectively.
Thus, for $C_\pm=C_{\mathrm{nt}}^{I=\pm1}$,
\begin{equation}\label{eq:descent-boundary-traces}
 \begin{aligned}
 \Tr(\Frob_p\mid W_4^+)&=\Tr(\Frob_p\mid C_+)-1,\\
 \Tr(\Frob_p\mid W_4^-)&=\Tr(\Frob_p\mid C_-)-\chi(p)
 \end{aligned}
 \qquad(p\ge5).
\end{equation}
In particular the second boundary contribution is $+1$ at inert primes.

\subsection{Alternating pairing and reciprocity}

\begin{proposition}\label{prop:descent-pairing}
Each $W_4^\pm$ has a perfect alternating $G_\Q$-equivariant pairing
with values in $\Q_2(-5)$. Its similitude character is $\tau^5$ and
its determinant is $\tau^{10}$.
\end{proposition}

\begin{proof}
The elliptic cup pairing $\langle\ ,\ \rangle$ on $R^1\pi_*\Q_2$
is alternating, perfect, and valued in $\Q_2(-1)$. On $\mathcal A_0$ put
\[
 B(v_1v_2v_3v_4,w_1w_2w_3w_4)
   =\sum_{\sigma\in S_4}\prod_{j=1}^4\langle v_j,w_{\sigma(j)}\rangle.
\]
This is symmetric with values in $\Q_2(-4)$. In a symplectic basis
$e,f$, its monomial matrix is antidiagonal with entries
$(-1)^j j!(4-j)!\langle e,f\rangle^4$, so it is perfect over $\Q_2$.
The degree-one elliptic isomorphisms $\phi_u$ preserve cup product;
hence $\theta$ preserves $B$.

Poincare duality and the compact-support map give a perfect parabolic
pairing $\Omega:H\times H\to\Q_2(-5)$
\cite[Theorem 24.1]{Milne}; see also the support and trace formalism
of \cite{SGA4}. To see perfection on the image, the two maps
$H_c^1\to H^1$ for dual coefficient systems are adjoint up to the
degree-one cup sign, and their kernels are the annihilators of the opposite
images. The degree-one cup sign and symmetry of $B$ make $\Omega$
alternating. Both $A$ and $I$ preserve it: their coefficient maps
are isometries and their curve automorphisms have degree one, so
preserve the trace map on $H_c^2$.

The adjoint of $A$ is $A^{-1}$; thus $e_{\mathrm{nt}}$ is self-adjoint
and its image has nondegenerate restricted pairing. Since $I$ is an
involutive isometry, its two eigenspaces are orthogonal and each is
nondegenerate. The restricted forms give the asserted equivariant
pairings. Their multiplier is the character of $\Q_2(-5)$, namely
$\tau^5$. The square of a nondegenerate alternating two-form in
dimension four is a volume form, yielding $\det W_4^\pm=\tau^{10}$.
\end{proof}

\begin{corollary}\label{cor:descent-reciprocity}
For $p\ge5$, the characteristic polynomial of geometric Frobenius
on either constructed representation is
\[
 \det(X-\Frob_p\mid W_4^\pm)
   =X^4-c_{1,\pm}(p)X^3+c_{2,\pm}(p)X^2
       -p^5c_{1,\pm}(p)X+p^{10}.
\]
The analogous formula for $V_4$ at a good prime of $K$ uses its norm
$q$ in place of $p$.
\end{corollary}

\begin{proof}
The pairing identity $F^t\Omega F=p^5\Omega$ pairs eigenvalues with
their quotients $p^5/\alpha$, and the determinant is $p^{10}$.
These identities give the displayed coefficients. Restriction to
$G_K$ and scalar extension transfer the pairing to $V_4$.
\end{proof}

\begin{remark}
This argument establishes reciprocity without a purity premise. The
coefficients of the displayed polynomials lie in $\Q_2$; a rational
compatible-system assertion, Hodge numbers and bad-prime factors
require their own arguments. Good-prime purity is proved next by a
separate application of Deligne's weight theorem; the rational motivic
realization and Hodge numbers are established in Section~\ref{sec:scholl},
and the local factors in Sections~\ref{sec:local-three} and \ref{sec:local-two}.
\end{remark}

% Insert after the descent/pairing results; all referenced results then precede this.
% Standard-theorem consequence, not a new general purity theorem.
\subsection{Good-prime Frobenius purity}

\begin{proposition}[Good-prime purity]\label{prop:good-prime-purity}
For $r=2,4$, the representation $H_r$ is pure of weight $r+1$ at
every prime $\mathfrak p\nmid6$ of $K$. More explicitly, if
$q=N\mathfrak p$ and $\alpha$ is a geometric Frobenius eigenvalue
on $H_r$, then $\alpha$ is algebraic over $\Q$ and
$|\iota(\alpha)|=q^{(r+1)/2}$ for every complex embedding of its
number field.
\end{proposition}

\begin{proof}
Fix $\mathfrak p\nmid6$, and work on
$U_0=\PP^1_{\F_q}\setminus\{0,1,\infty\}$. This is a smooth
curve of finite type over $\F_q$, and the coefficient prime $2$
is different from its characteristic. The specialized Legendre
family is smooth and proper over $U_0$, so
$\mathcal V_0=R^1\pi_*E$ is lisse. At any closed point $x$ of
degree $d$, its geometric Frobenius eigenvalues are the two
roots of $X^2-a_{q^d}(x)X+q^d$; the Weil theorem for elliptic
curves gives algebraic eigenvalues of absolute value $q^{d/2}$
under every complex embedding. Hence $\mathcal V_0$ is
pointwise pure of weight one.

The cover $z^3=\lambda$ is finite etale on $U_0$. Its chosen
rank-one character summand is lisse and has finite monodromy,
with Frobenius eigenvalues in $\mu_3$ at every closed point;
it is therefore pointwise pure of weight zero. Thus
$\mathcal G_{r,0}=\Sym^r\mathcal V_0\otimes\mathcal L_{\rho,0}$
is lisse and pure of weight $r$. The claims hold after extending
coefficients to $\overline{\Q}_2$ and for every embedding used
to define weights. Finite coefficient extension and passage to
the specified direct summand preserve these algebraicity and absolute-value assertions.

Deligne's \emph{Weil II}, Corollary~3.3.6 \cite{DeligneWII}, states that for a
smooth finite-type scheme and a lisse pure sheaf of weight $w$,
the image of $H_c^i\to H^i$ is pure of weight $w+i$.
Applying this result with $i=1$ proves that the special-fibre
parabolic group is pure of weight $r+1$. Equivalently, on this
curve one can apply Theorem~3.2.3 to
$H^1(\PP^1_{\overline{\F}_q},j_*\mathcal G_{r,0})$.
The specialization isomorphism of Section~\ref{sec:geometry}
is Frobenius equivariant and identifies this image with $H_r$,
so the asserted arithmetic purity follows. Neither argument
requires irreducibility or semisimplicity of $H_r$.
\end{proof}

\begin{corollary}\label{cor:descent-good-prime-purity}
Both $W_4^+$ and $W_4^-$ are pure of weight five at every rational
prime $p\ge5$: each geometric Frobenius eigenvalue is algebraic
and all its complex conjugates have absolute value $p^{5/2}$.
\end{corollary}

\begin{proof}
By Theorem~\ref{thm:descent-rational}, the restrictions of
$W_4^\pm$ to $G_K$, after extending coefficients, are $H_4$.
At a split prime $p$ the purity just proved applies to
$\Frob_p$ itself. At an inert prime, $\Frob_p^2$ restricts to
the geometric Frobenius at the degree-two prime of $K$.
Thus, if $\beta$ is an eigenvalue of $\Frob_p$ on either
$W_4^\pm$, then $\beta^2$ is algebraic and every complex
conjugate has absolute value $(p^2)^{5/2}=p^5$.
It follows that $\beta$ is algebraic and all its complex
conjugates have absolute value $p^{5/2}$, as required.
\end{proof}

\begin{remark}
This proposition concerns good-prime Frobenius weights of the
specified $2$-adic representations. It does not by itself
identify their Hodge numbers, construct rational Scholl
projectors or a compatible system, determine bad-prime Euler
factors or conductors, or prove automorphy or analytic
continuation. The separate construction in Section~\ref{sec:scholl}
supplies the specified rational Chow projectors and Hodge numbers.
Sections~\ref{sec:local-three}--\ref{sec:conductors} independently determine
the bad-prime Euler factors and conductors.
\end{remark}


\subsection{Projector compatibility on the smooth locus}
\begin{lemma}\label{lem:open-projectors}
Let $\mathcal F=\mathcal E\times_U T$ and
$h:Z=\mathcal F^{\times_T4}\to T$. On relative cohomology set
\[
 e_{\rm odd}=\prod_{j=1}^4\frac{1-[-1]_j^*}{2},\qquad
 e_{\rm alt}=\frac1{24}\sum_{\sigma\in S_4}\operatorname{sgn}(\sigma)\sigma^*,
 \qquad e_4=e_{\rm odd}e_{\rm alt}.
\]
These are rational commuting idempotents and
\[
 e_4R^j h_*\Q_2=0\quad(j\ne4),\qquad
 e_4R^4h_*\Q_2\simeq\Sym^4(R^1\pi_*\Q_2)|_T.
\]
The diagonal map induced by $\phi_u$ on $u^2=\lambda$ commutes with
$e_4$ and descends on this summand to the literal coefficient map
$\theta$ of Lemma~\ref{lem:descent-involution}. On parabolic cohomology,
$P_\pm=e_{\rm nt}(1\pm I)/2$ are orthogonal rational idempotents
with images $W_4^\pm$ and sum $e_{\rm nt}$.
\end{lemma}
\begin{proof}
Elliptic inversion acts by $-1$ on $R^1$ and by $+1$ on $R^0,R^2$.
Relative Kunneth therefore shows that $e_{\rm odd}$ selects exactly
the fourfold tensor of $R^1$, in degree four. Permuting these four
odd-degree factors introduces the Koszul sign. Consequently the signed
permutation average becomes the ordinary symmetric tensor projector
under Kunneth. This explains the sign used in Scholl's projector
\cite[Section~1.1.2]{SchollMotives}.

The diagonal elliptic isomorphism commutes with individual inversions
and factor permutations. Changing $u$ to $-u$ composes it with
simultaneous inversion on the four factors, acting by $(-1)^4=1$
on the selected summand. Its descended map is therefore exactly the
fourth symmetric power used to define $\theta$; the cocycle follows
from $\phi_{1/u}\phi_u=1$. The deck map commutes with $e_4$.
Since $IAI^{-1}=A^{-1}$, $I$ commutes with $e_{\rm nt}$, and the
last assertions follow by multiplying the displayed idempotents.
\end{proof}

This lemma establishes the smooth-family compatibility. We now extend it
to rational Chow projectors on a specified compactification, including the
auxiliary-level translations and the parabolic boundary comparison.

\section{Rational Scholl realization and Hodge numbers}\label{sec:scholl}
Write $\zeta=\zeta_K$. In this section $\pi$ also denotes the pullback
of the Legendre family to the displayed open base.
For any prime $\ell$, put
\[
 H_\ell=H^1_!(T_{\overline\Q},\Sym^4 R^1\pi_*\Q_\ell),\qquad
 W_{4,\ell}^{\pm}=(e_{\nt}H_\ell)^{I=\pm1},
\]
using the same literal deck and inversion maps as in
Section~\ref{sec:descent}. At $\ell=2$ these are the already constructed
$W_4^\pm$, without a change of sign or a Tate twist.

\begin{theorem}[Rational Scholl realization]\label{thm:scholl-realization}
There are a smooth projective, possibly disconnected fivefold $X/\Q$
and orthogonal idempotents
$p_\pm\in\CH^5(X\times_\Q X)\otimes\Q$ such that for every prime $\ell$
\[
 p_\pm H^j_{\et}(X_{\overline\Q},\Q_\ell)=0\quad(j\ne5),\qquad
 p_\pm H^5_{\et}(X_{\overline\Q},\Q_\ell)\simeq W_{4,\ell}^\pm.
\]
The corresponding Betti realization is compatible with Betti--etale
comparison and has $h^{5,0}=h^{0,5}=2$, with all other Hodge numbers zero.
Thus $(X,p_\pm,0)$ are rational Chow motives realizing the normalized
descents. The construction does not require irreducibility.
\end{theorem}
The following construction proves the theorem. Its use of Scholl's
standard projector theorem is separated from the explicit cover, rational
descent and sign computations for the present family.

\subsection{A semistable auxiliary surface}
Use the finite etale double cover
\[
S=\PP^1_s\setminus\{0,\infty,\mu_6\}\longrightarrow T,
\qquad z=s^2,\qquad \lambda=s^6.
\]
On $S$ the elliptic family is $E_s:y^2=x(x-1)(x-s^6)$. Its Weierstrass
invariants are
\[
\Delta=16s^{12}(1-s^6)^2,\qquad c_4=16(1-s^6+s^{12}).
\]
In characteristic zero, these give an $I_{12}$ fibre at $s=0$ and an $I_2$
fibre at each of the six roots of $s^6=1$. At infinity put
$t=1/s$, $X=t^6x$, $Y=t^9y$; the equation becomes
$Y^2=X(X-1)(X-t^6)$, giving another $I_{12}$ fibre.
The minimal regular model is therefore semistable, with no other singular
fibres. Its nonconstant $j$-invariant makes the family strictly nonconstant.

Now take the \emph{full frame torsor}
\[
S'=\operatorname{Isom}_S((\Z/12\Z)^2,E[12])\longrightarrow S.
\]
It includes all frames, not just one fixed Weil-pairing component and not
frames modulo their sign. It is a finite etale torsor under
$G=\GL_2(\Z/12\Z)$, defined over $\Q$, although it can be disconnected.
Let $\overline S'$ be its smooth projective normalization over $\PP^1_s$,
and let $\overline E'\to\overline S'$ be the minimal regular elliptic model.
Write $\pi':\overline E'|_{S'}\to S'$ for its smooth restriction, and
$\omega'$ for the Hodge line bundle on $\overline S'$.

\textbf{Why level 12 meets every component.} Around a geometric $I_m$ cusp,
the local monodromy on a full level-12 frame is conjugate to
\[
\begin{pmatrix}1&m\\0&1\end{pmatrix}\pmod {12}.
\]
Its order is $e=12/\gcd(12,m)$, the geometric ramification index of $S'$.
After base change and minimal regularization the fibre is $I_{me}$.
Here $m=12$ or $2$, so $me=12$ in both cases. In the Tate description,
the points represented by $q^{a/12}\zeta_{12}^b$ meet all twelve components
of an $I_{12}$ fibre. The tautological frame makes these points sections.
Thus the constant group $V=(\Z/12\Z)^2$ extends into the smooth locus and
meets every geometric fibre component. Any splitting needed in this local
description can be checked after a finite extension; the torsor itself is rational.

These facts verify the two geometric hypotheses of Scholl \cite[Section~2.2]{SchollVanishing}.
On each geometric component one can first enlarge its number field of
constants, apply that result, and then retain the entire finite disjoint union.
No connected-component choice is needed for the rational construction.

Form $(\overline E')^4_{\overline S'}$ and apply Scholl's canonical resolution;
call the result $X$. It has dimension five. The resolution blows up the
intrinsic strata determined by how many elliptic coordinates are nodes
\cite[Section~3.0]{SchollMotives}. The nodal local models suffice \cite[Section~2.10]{SchollVanishing}. Consequently every
family isomorphism considered below preserves its centers and lifts to $X$.


\subsection{Literal automorphisms and the sign relation}
On the family over $S$, define
\[
\begin{aligned}
a(s,x,y)&=(\zeta^2s,x,y),& d(s,x,y)&=(-s,x,y),\\
b(s,x,y)&=(s^{-1},x/s^6,y/s^9),& c(s,x,y)&=(s,x,-y).
\end{aligned}
\]
Use the same maps diagonally on four fibre coordinates. The order-three
choice $a(s)=\zeta^2s$ covers the required $z\mapsto\zeta z$.
For $b$, the chosen square root is exactly $u=s^3$.

A point of the frame cover carries a frame $v:V\to E_s[12]$. Lift $a$ and
$d$ by retaining this frame under the identity elliptic identification, and
lift $b$ by transporting it to $\phi_{s^3}\circ v$. For $h\in G$, let $h$
act by reparameterizing the frame while leaving elliptic points fixed.
Write $H_-$ for $h=-\mathrm{Id}$ and $C$ for simultaneous inversion of all
four fibre points, with the frame fixed. Upper-case $A,B,D$ denote the
pullback correspondences associated to the lower-case maps; $g^*$ is
represented by the transpose of its graph, with the pullback composition
convention used throughout.

The coordinate formulas give the following identities, already on the open
family with frames:
\[
\begin{gathered}
A^3=1,\qquad B^2=D^2=1,\qquad AD=DA,\qquad BAB=A^{-1},\\
\boxed{BDB=H_-CD.}
\end{gathered}
\]
For the last formula, applying $b,d,b$ sends
$(s,v,P_1,\ldots,P_4)$ to $(-s,-v,-P_1,\ldots,-P_4)$.
Omitting the frame sign or asserting $BD=DB$ on the whole space would be an
error. Both signs will disappear on the precise summand used below.

\textbf{Extension across the cusps.} An isomorphism of the generic elliptic
curve extends to its minimal regular model by uniqueness \cite{StacksMinimal}. This applies
also to the translated sections and to base automorphisms permuting cusps.
The local extensions agree on the dense open set and hence glue. The maps
then preserve the nodal strata used in the canonical fibre-product resolution,
so all act as genuine automorphisms of $X$, and their displayed relations
remain exact there. This uses that specific resolution, not an arbitrary one.

Frame reparameterizations commute with $A,B,D$. They normalize the universal
translation subgroup. The maps $A,B,D$ transport the frame and fibre points
together, so they normalize the translations, individual inversions, and
permutations, preserving the character used in the following construction.

All relevant automorphisms generate a finite group: the translation, inversion,
permutation and frame groups are finite, and the displayed relations give a
finite quotient generated by $A,B,D$. We may therefore use their rational
group algebra and its graph correspondences without a conjectural lifting of
cohomological idempotents.


\subsection{Exact rational Chow projectors}
Let
\[
\Gamma=(V\rtimes\{\pm1\})^4\rtimes\mathfrak S_4.
\]
Its character $\epsilon$ is trivial on translations, is the product of the
four inversion signs, and is $\sgn$ on $\mathfrak S_4$. Define
\[
\Pi_\epsilon=\frac1{|\Gamma|}\sum_{\gamma\in\Gamma}
\epsilon(\gamma)\gamma^*,\qquad
e_G=\frac1{|G|}\sum_{h\in G}h^*,\qquad q=e_G\Pi_\epsilon.
\]
The signed permutation is essential: the Koszul sign from exchanging degree-one
cohomology factors makes this select $\Sym^4 R^1$, rather than its alternating
power. The translation average supplies the component condition at the boundary.

The group $G$ normalizes $\Gamma$ and preserves $\epsilon$. Thus $e_G$ and
$\Pi_\epsilon$ commute, and $q^2=q$. Moreover $q$ commutes with $A,B,D$.
There are exact identities in the correspondence algebra
\[
H_-q=q,\qquad Cq=\epsilon(C)q=(-1)^4q=q.
\]
Hence $BDBq=Dq$, and $BDq=DBq$. Define
\[
d_+=\frac{1+D}{2},\qquad e_\nt=\frac{2-A-A^2}{3},\qquad
b_\pm=\frac{1\pm B}{2}
\]
Their restrictions $qd_+$, $qe_\nt$, and $qb_\pm$ are commuting idempotents
in $q\Q[\mathcal G]q$, whose identity is $q$ and where $\mathcal G$ is the
finite group just constructed. Indeed $B$ conjugates $A$ to $A^{-1}$, so it
commutes with $e_\nt$; $D$ commutes with $A$; and the remaining commutator
has just been removed by $q$. Consequently
\[
\boxed{p_\pm=q\,\frac{1+D}{2}\,\frac{2-A-A^2}{3}\,\frac{1\pm B}{2}}
\]
satisfy $p_\pm^2=p_\pm$, $p_+p_-=p_-p_+=0$, and
$p_++p_-=qd_+e_\nt$. These are equalities modulo rational equivalence,
because they already hold for the finite sums of automorphism correspondences.

\textbf{Descent to $\Q$.} The frame torsor, $X$, $B,D$, and $q$ are defined
over $\Q$. Although $A$ is defined over $\Q(\zeta)$, its Galois conjugate
is $A^2$. Therefore $A+A^2$, $e_\nt$, and $p_\pm$ are rational algebraic
correspondences. More formally, restriction and corestriction of rational Chow
cycles descend these invariant sums; restriction is injective after tensoring
with $\Q$. Thus the idempotent equalities hold over $\Q$, not only over
$\Q(\zeta)$.

No standard conjecture is used here. This is a decomposition by actual finite
automorphisms; it does not assert a Hecke-eigenform Chow decomposition or any
irreducibility of the resulting rank-four realizations by itself.
Section~\ref{sec:irreducibility} gives that separate argument for $\ell\ne3$.


\subsection{Parabolic realization and comparison}
Scholl's finite-group construction \cite[Theorem~1.2.1]{SchollMotives}, in the semistable
setting justified in \cite[Sections~2.7--2.12]{SchollVanishing}, gives for Betti and etale cohomology
\[
\Pi_\epsilon H^j(X)=0\quad(j\ne5),\qquad
\Pi_\epsilon H^5(X)\simeq
H^1_!(S',\Sym^4 R^1\pi'_*\Q).
\]
Here and below the coefficients on the etale side are $\Q_\ell$ and the
base is geometrically extended. The identification is induced by restriction
from the compact variety to the open fibre product, then relative Kunneth
and Leray. It is compatible with every family automorphism above.

There is a useful descent subtlety. Some choices identifying the cusp terms
in Scholl's boundary ladders become equivariant only after a finite extension.
We do not use those chosen boundary identifications. The proper-to-open
restriction, and the assertion that its image is the parabolic image, are
natural. Their injectivity and image equality may be verified after a finite
extension, and then hold for the same maps over $\Q$. The resulting
identification is fully $G_\Q$-equivariant. Nor is it necessary, or correct,
to say that every boundary cohomology group is killed by $\Pi_\epsilon$.

For a finite etale torsor with a finite group, pullback and trace identify
rational cohomology downstairs with invariants upstairs. The statement holds
for both ordinary and compactly supported cohomology, commutes with the
natural map between them, and therefore holds for its image. First take $G$
invariants and then $D$ invariants. This gives
\[
\begin{aligned}
qH^5(X)&\simeq H^1_!(S,\Sym^4 R^1\pi_{S*}\Q),\\
qd_+H^5(X)&\simeq H^1_!(T,\Sym^4 R^1\pi_*\Q).
\end{aligned}
\]
On the latter space $A$ is exactly the manuscript's cubic deck pullback.
The map $B$ uses $u=s^3$ in $\phi_u$, so its induced map is exactly $I$.
The discrepancy under $s\mapsto-s$ is simultaneous elliptic inversion,
already killed on $\Sym^4$. Applying $e_\nt$ and $b_\pm$ proves
\[
p_\pm H^5(X)\simeq(e_\nt H^1_!(T,\Sym^4 R^1\pi_*\Q))^{I=\pm1}.
\]
This proves the theorem's etale assertion for every $\ell$, in particular
the required equality with the existing $2$-adic $W_4^\pm$.

All maps used are pullback, trace, algebraic automorphisms, and natural
cohomological maps. Betti--etale comparison is functorial for these operations,
so the same formula identifies the rational Betti realization. Since it is a
projected subspace of $H^5$ of a smooth projective variety, its Hodge structure
is pure of weight five. Its precise Hodge numbers are computed next.


\subsection{Hodge numbers and an independent sign check}
The Legendre classifying map to the moduli \emph{stack} of elliptic curves is
etale: it factors as a $C_2$-torsor over the full-level-2 stack followed by an
$\mathfrak S_3$-torsor \cite[Lemma~4.3 and Proposition~4.5]{AntieauMeier}. This includes the
elliptic points where the coarse $j$-map ramifies. The maps $s\mapsto s^6$
and $S'\to S$ are etale on the open bases. Apply Scholl \cite[Section~2.13]{SchollVanishing} on $S'$,
where all the section hypotheses were verified, then take $G$ invariants.
The parabolic Hodge types are only $(5,0)$ and $(0,5)$, with
\[
H^{5,0}_!(S,\Sym^4 R^1\pi_{S*})
\simeq H^0(\PP^1_s,\omega^4\otimes\Omega^1_{\PP^1_s}),
\]
where $\omega$ is the Hodge line bundle on $\PP^1_s$. To justify this section
formula after descent, write $r:\overline S'\to\PP^1_s$. Semistable base
change gives $\omega'=r^*\omega$, and
$(r_*\Omega^1_{\overline S'})^G=\Omega^1_{\PP^1_s}$; locally this is the
identity $ds=e u^{e-1}du$ for $s=u^e$ in characteristic zero. The projection
formula then gives the displayed equality of holomorphic spaces.

The singular-fibre multiplicities sum to $12+12+6\cdot2=36$, giving
$\deg\omega=3$ and $\omega^4\otimes\Omega^1\simeq\mathcal O(10)$.
The same degree follows directly from the two minimal charts. Write
$\omega_s=dx/y$ and $\omega_t=dX/Y$ for relative invariant differentials.
On their overlap $\omega_s=t^3\omega_t$. Thus
\[
\beta_j=s^j\,ds\,\omega_s^4
       =-t^{10-j}\,dt\,\omega_t^4,\qquad 0\le j\le10,
\]
are a basis of holomorphic sections; the exponents show regularity at infinity.
The fourth power denotes the symmetric coefficient power, not a fourth wedge
of a one-form on the same elliptic curve.

Since $D^*\beta_j=(-1)^{j+1}\beta_j$, its invariant part has $j=1,3,5,7,9$.
With $z=s^2$ and $\omega_z=dx/y$ on the open family over $T$, let
\[
\alpha_k=z^k\,dz\,\omega_z^4\quad(0\le k\le4);
\qquad (s\mapsto s^2)^*\alpha_k=2\beta_{2k+1}.
\]
These give the five-dimensional $(5,0)$ part on $T$. The literal maps yield
\[
A\alpha_k=\zeta^{k+1}\alpha_k,\qquad I\alpha_k=-\alpha_{4-k}.
\]
For the second identity, use the target coordinate $s'$: $B^*\omega_{s'}=s^3\omega_s$
and $B^*(ds')=d(s^{-1})=-s^{-2}ds$, hence $B^*\beta_j=-\beta_{10-j}$.
This fixes the sign without any appeal to numerical traces.

The projector $e_\nt$ removes precisely $\alpha_2$. On the remaining four
holomorphic dimensions the two eigenspaces are
\[
\begin{array}{c|c}
I=+1&\alpha_0-\alpha_4,\quad\alpha_1-\alpha_3\\[2pt]
I=-1&\alpha_0+\alpha_4,\quad\alpha_1+\alpha_3.
\end{array}
\]
Each has dimension two. Rationality of $p_\pm$ and complex conjugation give
two antiholomorphic dimensions as well. Endpoint Hodge type excludes all
other $(p,q)$. Each realization therefore has rank four and the asserted
Hodge numbers, independently confirming the prior rank calculation.


The construction and the Hodge calculation complete the proof of
Theorem~\ref{thm:scholl-realization}.

\section{Inert Frobenius traces and good Euler factors}\label{sec:inert}

Retain the Legendre family $\mathcal E_\lambda$, its cubic cover
$f:T\to U$, and the operators $A$ and $I$ from
Sections~\ref{sec:geometry} and~\ref{sec:descent}.
Thus $A$ is pullback by $z\mapsto\zeta_3z$, while $I$ is induced by
$z\mapsto z^{-1}$ and the fourth symmetric power of
\[
 \phi_u:\mathcal E_\lambda\longrightarrow\mathcal E_{1/\lambda},
 \qquad(t,y)\longmapsto(t/\lambda,y/u^3),\qquad u^2=\lambda.
\]
In particular, $I^2=1$, $IAI^{-1}=A^{-1}$, and $I$ commutes with
Galois. Write $H_{\nt}$ for the eight-dimensional nontrivial deck
summand and $W_4^\pm=H_{\nt}^{I=\pm1}$. Frobenius on cohomology
always means geometric Frobenius.

Fix a prime $p\geq5$ with $p\equiv2\pmod3$. Choose a nonsquare
$d\in\F_p^\times$, a root $r\in\F_{p^2}$ of $r^2=d$, and a
primitive cubic character $\rho:\F_{p^2}^\times\to\mu_3$.
For $x\in\F_p$ set
\begin{align*}
 \nu_x&=x+r,&
 E'_{x,d}&:\quad Y^2=X(X^2-2xX+x^2-d),\\
 a'_p(x)&=p+1-\#E'_{x,d}(\F_p),&
 P_4(a,p)&=a^4-3pa^2+p^2,
\end{align*}
and define
\begin{equation}\label{eq:inert-sum}
 S_p=\sum_{x\in\F_p}\rho(x+r)P_4(a'_p(x),p).
\end{equation}
Every fibre in this sum is smooth: its three roots $0,x+r,x-r$
are distinct. Character values can be regarded in $\Q(\zeta_3)$
and then embedded in the coefficient field $E=\Q_2(\zeta_3)$.

\begin{theorem}\label{thm:inert-trace}
The sum $S_p$ is an integer independent of $d$, the sign of $r$,
and the primitive cubic character. At every such prime,
\begin{equation}\label{eq:inert-traces}
 \Tr(\Frob_p\mid W_4^+)=-S_p-1,
 \qquad
 \Tr(\Frob_p\mid W_4^-)=S_p+1.
\end{equation}
\end{theorem}

\begin{proof}
We realize the operator $I\Frob_p$ by an actual variety and sheaf
over $\F_p$, so that the ordinary trace formula applies.
Let
\[
 C=\ker\!\left(N_{\F_{p^2}/\F_p}:
       \operatorname{Res}_{\F_{p^2}/\F_p}\mathbb G_m
       \longrightarrow\mathbb G_m\right).
\]
This torus has equation $a^2-db^2=1$. Over $\F_{p^2}$, the
coordinate $z=a+rb$ identifies $C$ with $\mathbb G_m$.
Since $r^p=-r$, arithmetic Frobenius on geometric points is
\begin{equation}\label{eq:torus-point-frobenius}
 z\longmapsto z^{-p}.
\end{equation}
Indeed $(a-rb)^p=a^p+rb^p$. Put $U'=C\setminus\{1\}$ and
$T'=C\setminus C[3]$. The cube map $h:T'\to U'$ is finite
\'{e}tale of degree three. Its deck group $C[3]$ is constant over
$\F_p$, because $\zeta^{-p}=\zeta$ for $\zeta\in\mu_3$.
Over $\F_{p^2}$ the map is the original cover $z^3=\lambda$.
The parameter
\begin{equation}\label{eq:torus-parameter}
 \lambda=\frac{x+r}{x-r},\qquad
 x=r\frac{\lambda+1}{\lambda-1}
\end{equation}
identifies $U'$ with $\mathbb A^1_{\F_p}\setminus V(x^2-d)$.
Thus $U'(\F_p)=\F_p$ in this parameter. The removed point
$\lambda=1$ corresponds to $x=\infty$; the other geometric
punctures are $x=r,-r$.

Let $\mathcal K_{\nt}$ be the rank-two augmentation summand of
$h_*\Q_2$. The elliptic curves $E'_{x,d}$ form an elliptic scheme
$\pi':\mathcal E'\to U'$, and we use the lisse sheaf
\begin{equation}\label{eq:inert-twisted-sheaf}
 \mathcal B=\Sym^4(R^1\pi'_*\Q_2)\otimes\mathcal K_{\nt}.
\end{equation}
We now check the coefficient descent on the whole family.
Over $\F_{p^2}$ write $a=x-r$, $\nu=x+r$, so that
$\lambda=\nu/a$ and $\bar a=\nu$. On the finite \'{e}tale
cover $s^2=a$ there is an elliptic isomorphism
\[
 \alpha:\mathcal E'\longrightarrow\mathcal E_\lambda,
 \qquad (X,Y)\longmapsto(X/a,Y/s^3).
\]
Replacing $s$ by $-s$ composes this map with $[-1]$; hence its
fourth symmetric power descends to $U'_{\F_{p^2}}$.
On a common square-root cover choose $\bar s$ with
$\bar s^2=\bar a=\nu$ and put $u=\bar s/s$. Then $u^2=\lambda$
and the maps of families satisfy
\begin{equation}\label{eq:inert-family-cocycle}
 \phi_u\circ\alpha(X,Y)
  =(X/\nu,Y/\bar s^3)=\bar\alpha(X,Y).
\end{equation}
This is exactly the compatibility of the induced pullbacks with
the coefficient map defining $I$. All square-root choices disappear
on $\Sym^4$. Effective \'{e}tale descent, applied at finite
coefficient levels and then to their compatible inverse system,
identifies \eqref{eq:inert-twisted-sheaf} with the inversion twist
of the original coefficient system \cite{StacksDescent}.
Its geometric Frobenius on cohomology is $IF$, where
$F=\Frob_p$ on the original cohomology: the arithmetic descent
uses the involution $I$, and $I^{-1}=I$ and $IF=FI$.
In particular, at a rational point the elliptic coefficient trace is
exactly $P_4(a'_p(x),p)$. This follows from the eigenvalue sum
$\alpha^4+\alpha^3\beta+\alpha^2\beta^2+\alpha\beta^3+\beta^4$
with $\alpha+\beta=a'_p(x)$ and $\alpha\beta=p$; no additional
scalar occurs.

For the Kummer coefficient, fix $x\in\F_p$, choose $v^3=\nu_x$,
and take the lift $z=v^{1-p}$ of $\lambda$. Then
$z^3=\nu_x^{1-p}=\nu_x/\bar\nu_x$. The arithmetic point map
\eqref{eq:torus-point-frobenius} takes $z$ to $\kappa z$, where
\begin{equation}\label{eq:inert-kummer-weight}
 \kappa=z^{-p-1}=v^{p^2-1}
        =\nu_x^{(p^2-1)/3}\in\mu_3.
\end{equation}
The result is independent of the cube root $v$.
For the pullback eigenspace of functions specified by
$g(\zeta_3z)=\zeta_E g(z)$, geometric Frobenius acts by pullback
under this arithmetic point permutation and has eigenvalue
$\iota(\kappa)$, where $\iota(\zeta_3)=\zeta_E$. Thus we may take
$\rho(\nu)=\iota(\nu^{(p^2-1)/3})$; the other nontrivial line
has inverse eigenvalue. The local trace of $\mathcal B$ is therefore
\begin{equation}\label{eq:inert-local-trace}
 \bigl(\rho(x+r)+\rho(x+r)^{-1}\bigr)P_4(a'_p(x),p).
\end{equation}
This computation distinguishes arithmetic action on points from
geometric action on cohomology.

Put $\mathcal A_0=\Sym^4R^1\pi_*\Q_2$ and
$C_{\nt}=e_{\nt}H_c^1(T_{\overline{\F}_p},f^*\mathcal A_0)$,
where $e_{\nt}=(2-A-A^2)/3$.
Finite \'{e}tale pushforward identifies this with the corresponding
compactly supported cohomology on the base \cite{StacksFinite}.
The geometric groups in degrees zero and two vanish by
Section~\ref{sec:geometry}, and twisting preserves those vanishings.
The trace formula \cite{StacksTrace} and
\eqref{eq:inert-local-trace} give
\[
 \Tr(IF\mid C_{\nt})=-(S_p+\bar S_p).
\]
The boundary calculation of Section~\ref{sec:geometry} gives
\[
 0\longrightarrow B_{\nt}\longrightarrow C_{\nt}
   \longrightarrow H_{\nt}\longrightarrow0,
 \qquad B_{\nt}=\Q_2[\mu_3]_{\nt}.
\]
The coefficient line at $\lambda=1$ has Frobenius eigenvalue $1$,
and the coefficient action of $I$ there is $+1$.
Both $F$ and $I$ permute $\mu_3$ by inversion. Hence $IF$ is the
identity on this two-dimensional boundary, and
\begin{equation}\label{eq:inert-parabolic-trace}
 \Tr(IF\mid H_{\nt})=-(S_p+\bar S_p)-2.
\end{equation}
Moreover, $FAF^{-1}=A^{-1}$, so $F$ exchanges the two nontrivial
$A$ eigenspaces and $\Tr(F\mid H_{\nt})=0$.
The projectors $(1\pm I)/2$ now give
\[
 \Tr(F\mid W_4^+)=-\frac{S_p+\bar S_p}{2}-1,
 \qquad
 \Tr(F\mid W_4^-)=\frac{S_p+\bar S_p}{2}+1.
\]

It remains to establish integrality and independence of choices.
The character $\rho$ is trivial on $\F_p^\times$, because
$3\nmid p-1$. Under $x\mapsto-x$ we have
$\nu_{-x}=-\bar\nu_x$ and hence
$\rho(\nu_{-x})=\rho(\nu_x)^{-1}$.
The curve $E'_{-x,d}$ is the quadratic twist by $-1$ of $E'_{x,d}$;
therefore $a'_p(-x)=\chi_p(-1)a'_p(x)$, where $\chi_p$ is the
quadratic character. Since $P_4$ is even, $S_p=\bar S_p$.
As $S_p\in\Z[\zeta_3]$, it follows that $S_p\in\Z$ and the
trace identities follow. If $d'=c^2d$ and $r'=cr$, relabel
$x'=cx$. The character weight is unchanged, and $E'_{cx,d'}$
is the quadratic twist by $c$ of $E'_{x,d}$, again invisible to
$P_4$. Every two nonsquares differ by such a square.
Changing the sign of $r$ or replacing $\rho$ by its inverse
conjugates the sum and therefore also leaves it unchanged.
\end{proof}

For the good Euler factors, let
\[
 v(q)=-M_{3,4}(q)+3qM_{3,2}(q)+q^2-1
\]
be the $V_4$ trace established in Section~\ref{sec:geometry}.
Whenever $q$ is a higher power of the residue field of a fixed
prime, the cubic character in these moments is its norm lift.

\begin{corollary}\label{cor:good-euler-factors}
For every rational prime $p\geq5$, put
\[
 c_1(p)=
 \begin{cases}
  v(p),&p\equiv1\pmod3,\\
  S_p+1,&p\equiv2\pmod3,
 \end{cases}
 \qquad
 c_2(p)=\frac{c_1(p)^2-v(p^2)}{2}.
\]
Then the good Euler polynomial of $W_4^-$ is
\begin{equation}\label{eq:good-euler-polynomial}
 \det(1-\Frob_pT\mid W_4^-)
 =1-c_1(p)T+c_2(p)T^2-p^5c_1(p)T^3+p^{10}T^4.
\end{equation}
\end{corollary}

\begin{proof}
At split primes, either extension restricts to $V_4$ and has trace
$v(p)$. At inert primes Theorem~\ref{thm:inert-trace} gives the
stated first trace. In both cases
$\Tr(\Frob_p^2\mid W_4^-)=v(p^2)$: at a split prime this uses
the norm-lifted character from the chosen residue field $\F_p$;
at an inert prime $\F_{p^2}$ is the initial residue field.
Newton's identity gives $c_2(p)$. The perfect alternating pairing
of Section~\ref{sec:descent}, with Frobenius multiplier $p^5$,
gives determinant $p^{10}$ and the coefficient $-p^5c_1(p)$
of $T^3$.
\end{proof}

Thus the exact finite-sum prescription with inert first trace $S_p+1$
is the good-prime Euler prescription of $W_4^-$. This statement
does not verify numerical implementations or floating-point tables.
The separate local proofs in Sections~\ref{sec:local-three} and
\ref{sec:local-two} supply the factors at $2$ and $3$. An identification
of completed analytic $L$-functions remains outside the established scope.

\section{Local inertia and the Euler factors at 3}\label{sec:local-three}
This section computes the arithmetic local representation of the literal
rational descents fixed in Section~\ref{sec:descent}. It combines an explicit
quotient with the local two-section vanishing-cycle calculation and a separate
Frobenius computation. The proof does not use a numerical functional-equation fit.

\subsection{Statement and normalization}
Let $\ell\ne3$, let
\[
 T=\PP^1_z\setminus\{0,\infty,\mu_3\},\qquad
 E_z:y^2=x(x-1)(x-z^3),\qquad
 H_\ell=H^1_!(T_{\overline\Q},\Sym^4 R^1\pi_*\Q_\ell).
\]
Here $H^1_!=\im(H_c^1\to H^1)$, $A$ is pullback by $z\mapsto\zeta_3z$,
and $e_\nt=(2-A-A^2)/3$. The literal inversion $I$ uses
\[
 z\mapsto z^{-1},\qquad (x,y)\mapsto(x/z^3,y/u^3),\quad u^2=z^3.
\]
The square-root sign disappears on the fourth symmetric power. Put
$W_{4,\ell}^{\pm}=(e_\nt H_\ell)^{I=\pm1}$, as in \cite{WalkeV2}.
We use geometric Frobenius, so $\Q_\ell(-1)$ has eigenvalue $3$ at $3$.
The standing descent identities are
\[
 \dim W_{4,\ell}^{\pm}=4,\qquad
 W_{4,\ell}^{-}\simeq W_{4,\ell}^{+}\otimes\chi_{-3}.
\]
They follow from the same finite-cover and involution construction for every
$\ell$, without requiring compatibility at a bad prime.

\begin{theorem}\label{p3:thm:local}
For every $\ell\ne3$, the local quantities of these specified representations are
\[
\begin{array}{c|cccc}
 &\dim(W_{4,\ell}^{\pm})^{I_3}&P_3^{\pm}(T)&\Sw_3(W_{4,\ell}^{\pm})&a_3(W_{4,\ell}^{\pm})\\\hline
 +&1&1-3T&0&3\\
 -&0&1&0&4
\end{array}
\]
where $P_3^{\pm}(T)=\det(1-\operatorname{Frob}_3T\mid
(W_{4,\ell}^{\pm})^{I_3})$.
On $W_{4,\ell}^{+}$, inertia is unipotent with a single Jordan block of length
four; on $W_{4,\ell}^{-}$ it is its ramified quadratic twist.
If $N$ is the nonzero tame monodromy logarithm and $F$ is geometric
Frobenius, then $FNF^{-1}=3^{-1}N$. On the successive one-dimensional
quotients of $0\subset\ker N\subset\ker N^2\subset\ker N^3\subset V$
for $V=W_{4,\ell}^+$, the Frobenius eigenvalues are $3,9,27,81$.
\end{theorem}

In particular, the original $2$-adic representation can be used directly at
$3$. The argument also works at $\ell=5$ and yields the same displayed data.
The dyadic calculation is treated separately below.

\subsection{An exact quotient, including the inversion sign}
Set $t=z+z^{-1}$ and consider
\begin{equation}\label{p3:eq:family}
 C_t:Y^2=X^3+(t-1)(t+2)X^2+(t+2)X.
\end{equation}
Write $a=(t-1)(t+2)$, $b=t+2$. Its invariants are
\begin{align*}
 c_4&=16(a^2-3b)=16(t+2)(t^3-3t-1),\\
 \Delta&=16b^2(a^2-4b)=16(t+2)^3(t-2)(t+1)^2,\\
 j&=\frac{256(t^3-3t-1)^3}{(t-2)(t+1)^2}.
\end{align*}
Over $\Q(z)$ put $d=-z^2/(z+1)$. Direct identities give
\[
 da=-(1+z^3),\qquad d^2b=z^3.
\]
Thus $E_z$ is precisely the quadratic twist of $C_t$ with coefficients
$da,d^2b$. Over $v^2=d$, the map
\[
 \psi_z:E_z\longrightarrow C_t,\qquad(x,y)\longmapsto(x/d,y/v^3)
\]
induces a canonical fourth-symmetric-power isomorphism. Its two choices
differ by $[-1]$, which acts trivially on that symmetric power.

Crucially, $d(1/z)=d(z)/z^3$. Choosing $v'=v/u$ for the inverted parameter gives
\[
 \psi_{1/z}\circ\phi_u(x,y)=(x/d,y/v^3)=\psi_z(x,y).
\]
This proves compatibility with the \emph{literal} $I$ in the manuscript.
The descent is the plus descent; it has no additional constant quadratic twist.

Let $V=\PP^1_t\setminus\{2,-1,-2,\infty\}$ and
$Z=\PP^1_z\setminus\{0,\infty,\mu_3,-1\}$.
Then $Z\to V$ is a finite etale double cover, and the displayed isomorphism
identifies its invariant coefficient system with
$\mathcal F=\Sym^4R^1(C/V)_*\Q_\ell$.
On the compact curves, finite pushforward followed by invariants identifies
the ordinary middle-extension sheaves. This can be checked stalkwise:
the stalk at a missing point is the invariant space of its punctured-trait
monodromy. Finite pushforward and rational finite-group invariants are exact.
Consequently
\begin{equation}\label{p3:eq:plus}
 H^1(\PP^1_{\bar\Q},j_*\mathcal F)\simeq H_\ell^{I=+1}.
\end{equation}
Adding back $z=-1$ on the original side makes no change to its parabolic
cohomology, since its original coefficient sheaf is smooth there.

\subsection{The rank and the extra branch point}
The singular fibres of \eqref{p3:eq:family} in characteristic zero are
\[
\begin{array}{c|rrrr}
 t&2&-1&-2&\infty\\\hline
 \text{type}&I_1&I_2&\mathrm{III}&I_6\\
 \dim\mathcal F^{I_t}&1&1&3&1
\end{array}
\]
At infinity use $r=1/t$, $X'=r^2X$, $Y'=r^3Y$. The resulting coefficients
are $1+r-2r^2$ and $r^3(1+2r)$; $c_4$ is a unit and $\Delta$ has order six.
At $-2$, the orders of $c_4,\Delta$ are $1,3$, giving type III.
Equivalently, the ramified preimage $z=-1$ has elliptic automorphism
$(x,y)\mapsto(-x,iy)$. Its $H^1$ eigenvalues are $i,-i$, so the eigenvalues
of its fourth symmetric power are $1,-1,1,-1,1$.
This explains both the three-dimensional stalk at $-2$ and why that point
must be included in the rank calculation.

The Legendre geometric monodromy is $\mathrm{SL}_2$; finite-index pullback
preserves its connected Zariski closure. The quotient identification therefore
shows that neither $\mathcal F$ nor its dual has a geometric global invariant.
The parabolic Euler-characteristic formula gives
\[
 \dim H^1(\PP^1,j_*\mathcal F)=(4-2)5-(1+1+3+1)=4.
\]
The already defined rank-four $W_{4,\ell}^{+}$ is a subspace of the right
side of \eqref{p3:eq:plus}; equality of dimensions proves
\begin{equation}\label{p3:eq:identification}
 W_{4,\ell}^{+}=H^1_!(V_{\bar\Q},\mathcal F)
\end{equation}
with its full arithmetic action. This step does not use Hodge numbers or a
numerical Frobenius comparison.

\subsection{The collision at 3 and its local cycles}
On $\mathcal X=\PP^1_{\Z_3}$ remove the four sections
$D=\{2,-1,-2,\infty\}$ and let $j:\mathcal U\hookrightarrow\mathcal X$.
The elliptic scheme \eqref{p3:eq:family} is smooth and proper over $\mathcal U$,
since its discriminant is a unit there. Thus $\mathcal F$ is lisse on this
entire open scheme for every $\ell\ne3$.

The sections $t=2,-1$ meet at the special point $x$ with $t=-1\in\F_3$.
Their intersection multiplicity is $v_3(2-(-1))=1$. The other two sections
specialize to $1,\infty$, separately. Near $x$ the invariant $c_4$ is a unit
and the discriminant is a unit times $(t-2)(t+1)^2$.
At $x$ the nodal cubic is
\[
 Y^2=X^3+X^2+X=X(X-1)^2.
\]
The node $(1,0)$ has rational tangents $Y=\pm(X-1)$: the degeneration is
split multiplicative. The two horizontal monodromies on $R^1$ have
multiplicities one and two in the same Tate unipotent direction.

We apply \emph{the local} theorem \cite[\S\S1.6--1.9]{SchollVanishing}, not its
global semistable-family Theorem 2.15. The present family has an additive
point at $-2$, so that global theorem's hypotheses should not be asserted.
For the local theorem, the base curve is smooth, the two sections have
intersection multiplicity one, the coefficient system is lisse off them,
and Tate uniformization makes its local monodromy factor through the two
Kummer characters with multiplicities $1,2$. These are exactly the required
local hypotheses.

Put $B=(R^1\Phi j_!\mathcal F)_x$. Scholl's calculation identifies the
inertia action on $B$ with the rank-five fourth symmetric power of a
nontrivial unipotent rank-two action. Its logarithm is proportional to
$(1+2)N_0=3N_0$, and $B$ is a single unipotent block of length five.
In particular wild inertia acts trivially.
The local boundary sequence \cite[\S1.7]{SchollVanishing} gives
\begin{equation}\label{p3:eq:ABC}
 0\longrightarrow A_x\longrightarrow B\longrightarrow C_x\longrightarrow0,
 \qquad \dim A_x=1,\quad \dim C_x=4.
\end{equation}
Here $A_x$ is the cokernel of the diagonal map from the common invariant
line to the two cusp invariant lines. It is the unique invariant line of $B$.
Thus $C_x$ is a single unipotent block of length four.

\subsection{The global-to-local surjection}
This paragraph supplies the necessary bridge from local cycles to the
specified global representation. It also avoids any assumption about the
absence of other vanishing cycles at the additive point. In particular,
we do not invoke the original deduction of Corollary 2.16 in Scholl's
paper; its published correction is included in the reference
\cite{SchollVanishing}.

On the special open curve $\mathcal U_{\bar\F_3}$,
$H_c^2(\mathcal U_{\bar\F_3},\mathcal F)=0$.
Indeed pull back by the separable double cover $t=z+z^{-1}$.
The exact quadratic-twist identities still hold in characteristic three.
After taking $\Sym^4$, the resulting system is the Legendre system at $z^3$.
The relative Frobenius isogeny from the Legendre curve at $z$ to that at $z^3$
induces an isomorphism on $\ell$-adic $H^1$ for $\ell\ne3$.
The Legendre geometric monodromy is $\mathrm{SL}_2$ in every odd
characteristic \cite[Chapter 15, opening discussion]{Katz}; hence its
fourth symmetric power and its dual have no global invariants. Pullback
detects invariant sections, and Poincare duality gives the stated vanishing.

Apply the nearby/vanishing-cycle triangle to $\mathcal G=j_!\mathcal F$
on the proper curve $\mathcal X$. Since $\mathcal G$ is lisse on the smooth
open, $R\Phi\mathcal G$ is supported on the finite special boundary.
Its degree-one hypercohomology is therefore the direct sum of its
degree-one stalks. The exact sequence contains
\[
 H_c^1(V_{\bar\Q_3},\mathcal F)\longrightarrow
 \bigoplus_{y\in D_{\bar\F_3}}(R^1\Phi\mathcal G)_y
 \longrightarrow H_c^2(\mathcal U_{\bar\F_3},\mathcal F)=0.
\]
Projection gives a surjection onto $B$. Naturality of the boundary maps
and \eqref{p3:eq:ABC} show that the boundary subspace of $H_c^1$ maps into
$A_x$. More explicitly, the two distant cusp stalks have zero contribution
at $x$, while the two colliding invariant lines map through
$(\Q_\ell\oplus\Q_\ell)/\operatorname{diag}\Q_\ell=A_x$ and onto it.
Thus there is an induced
surjection
\[
 H^1_!(V_{\bar\Q_3},\mathcal F)\longrightarrow C_x.
\]
Its source and target have dimension four, so it is an isomorphism.
All constructions are natural over $\Q_3$, and $x$ is rational over $\F_3$;
the isomorphism is equivariant for the full local Galois group.
Together with \eqref{p3:eq:identification}, this proves that $W_{4,\ell}^{+}$
has tame, unipotent inertia with a single length-four block.

\subsection{Arithmetic Frobenius: fixing the sign and the Tate twist}
Scholl's local theorem by itself specifies inertia, not the arithmetic
Frobenius eigenvalue. We give the extra calculation here.
Near the split Tate degeneration, $R^1$ has its natural filtration
\[
 0\longrightarrow\Q_\ell\longrightarrow R^1
 \longrightarrow\Q_\ell(-1)\longrightarrow0.
\]
The fourth symmetric power has five graded pieces
$\Q_\ell(-j)$, $0\le j\le4$.
They are unramified arithmetic coefficient lines extending over the local
base. The split-node tangent calculation removes any torus splitting
character. Alternatively a quadratic splitting character disappears on
the even fourth power.

For a constant unramified coefficient line $L$, the local vanishing cycles
of extension by zero at two colliding, individually rational sections satisfy
\[
 R^1\Phi(j_!L)_x=\operatorname{coker}(L\xrightarrow{\rm diag}L\oplus L)
 \simeq L,\qquad R^q\Phi(j_!L)_x=0\quad(q\ne1).
\]
This follows naturally from the constant-sheaf/closed-boundary exact
sequence and smooth base change. Because both sections are separately
rational, Frobenius does not exchange them. In particular no sign
character and no extra Tate twist enter this formula.

Apply the vanishing-cycle functor to the Tate filtration. The displayed
concentration in degree one makes successive sequences short exact at this
stalk. It yields an arithmetic filtration of $B$ with graded pieces
\[
 \Q_\ell,\ \Q_\ell(-1),\ \Q_\ell(-2),\ \Q_\ell(-3),\ \Q_\ell(-4).
\]
The boundary line $A_x$ is the first piece: it is the difference of the
two split cusp invariant lines. Hence $C_x$ has graded pieces with
$j=1,2,3,4$. Its first one-dimensional filtered piece is inertia invariant;
the single length-four block has just one invariant line. Therefore
\[
 (W_{4,\ell}^{+})^{I_3}\simeq\Q_\ell(-1),\qquad P_3^+(T)=1-3T.
\]
This determines the positive sign arithmetically. It does not infer that
sign from purity, polarization or a numerical functional equation.

The same filtration determines the full ordered monodromy flag. Every
filtered piece is inertia-stable, and a single nilpotent Jordan block
has a unique invariant subspace of each dimension. Thus these pieces
are exactly $\ker N^r$, with graded Frobenius eigenvalues $3^r$ for
$r=1,2,3,4$. Geometric Frobenius acts on the tame inertia parameter by
$3^{-1}$, so taking logarithms gives $FNF^{-1}=3^{-1}N$.

\subsection{The minus descent and the conductor exponents}
The extension $\Q_3(\sqrt{-3})/\Q_3$ is ramified of degree two, hence tame.
Its character $\chi_{-3}$ is nontrivial on inertia and trivial on wild inertia.
On the plus space all inertia operators are unipotent. Choose an inertia
element on which $\chi_{-3}=-1$; its action on the minus space has all
eigenvalues $-1$, so it fixes no nonzero vector. Consequently
\[
 (W_{4,\ell}^{-})^{I_3}=0,\qquad P_3^-(T)=1.
\]
Both spaces are tame, so both Swan conductors vanish. The Artin conductor
formula \cite{Ulmer} is
\[
 a_3(V)=\dim V-\dim V^{I_3}+\Sw_3(V).
\]
It gives $a_3(W^+)=3$, $a_3(W^-)=4$, completing Theorem~\ref{p3:thm:local}.


\section{The local factor and conductor exponent at 2}\label{sec:local-two}
Let $\ell\ne2$ and retain the rational descents defined by the literal
cubic deck operation and inversion. On
$T=\PP^1_z\setminus\{0,\infty,\mu_3\}$ let
$f:\mathcal E\to T$ be the family $y^2=x(x-1)(x-z^3)$, and put
\[
 \mathcal H_\ell=H^1_!(T_{\overline\Q},\Sym^4R^1f_*\Q_\ell),
 \qquad W_{4,\ell}^{\pm}=(e_{\nt}\mathcal H_\ell)^{I=\pm1}.
\]
The following argument uses one arithmetic Galois cover for the complete
calculation. The semistable-model criterion supplies tameness; the main,
tail and node terms then determine the inertia invariants.

\begin{theorem}[Dyadic local data]\label{p2:thm:dyadic-local-data}
For every coefficient prime $\ell\ne2$,
\[
 \mathcal H_\ell^{I_2}=0,\qquad
 \Sw_2(\mathcal H_\ell)=0.
\]
Consequently both specified rank-four descents satisfy
\[
 d_{2,\ell}^{\pm}=0,\qquad P_{2,\ell}^{\pm}(T)=1,\qquad
 \Sw_2(W_{4,\ell}^{\pm})=0,\qquad a_2(W_{4,\ell}^{\pm})=4.
\]
\end{theorem}
The proof occupies the remainder of this section. It retains the
arithmetic action on the cover and on its coefficient system separately.
The identity $W^-_{4,\ell}=W^+_{4,\ell}\otimes\chi_{-3}$ is consistent
with the equal dyadic data: $\chi_{-3}$ is unramified at $2$ and has
Frobenius value $-1$ there.

\subsection{The exact Galois closure and its arithmetic action}
Fix an auxiliary prime $\ell\ne2$, put $K_0=\Q_2(\zeta)$, where
$\zeta^3=1\ne\zeta$, and let $L=K_0(\pi)$ with $\pi^3=2$.
Choose an arithmetic inertia lift $\sigma$ of the generator of
$\operatorname{Gal}(L/K_0)$ with $\sigma(\pi)=\zeta\pi$.
The following computations concern actual geometric deck transformations
and this arithmetic automorphism separately.

\subsubsection{The constant in the cubic closure}
Set $h=v^3$, $D=2h+1$, and $z=v(h+2)/D$.  The polynomial
\[
 f_z(X)=X^4-2zX^3+2X-z
\]
factors after adjoining $v$ as
\[
 f_z(X)=(X-v)\left(X^3-\frac{3v}{D}X^2
                 -\frac{3v^2}{D}X+\frac{h+2}{D}\right).
\]
Adjoin $\eta$ satisfying
\begin{equation}\label{p2:eq:closure-constant}
 \eta^3=4(v^3+1).
\end{equation}
The other three roots are exactly
\begin{equation}\label{p2:eq:closure-roots}
 v_i=\frac{v-\zeta^i\eta^2/2-\zeta^{-i}v^2\eta}{D},
 \qquad i=0,1,2.
\end{equation}
Indeed, after the substitution $X=Y+v/D$, the cubic has coefficients
\[
 p=-\frac{6v^2(h+1)}{D^2},\qquad
 q=\frac{2(h+1)(2h^2+h+1)}{D^3}.
\]
The two Cardano terms $u=-\eta^2/(2D)$ and $w=-v^2\eta/D$
satisfy $uw=-p/3$ and $u^3+w^3=-q$.
The discriminant is
\[
 \operatorname{disc}(f_z)=-432(z^3-1)^2
                         =-3\bigl(12(z^3-1)\bigr)^2.
\]
Consequently $K_0(v,\eta)$ is the splitting field over $K_0(z)$:
it contains the four roots and has degree $4\cdot3=12$, since
$v^3+1$ has simple zeros and is not a cube in $K_0(v)$.
Its group is $A_4$.  Over $\Q_2(z)$ the group is $S_4$ and the
constant extension is exactly $K_0$.  For example, the smooth cubic
has the rational point $(v,\eta)=(-1,0)$, so its function field is
geometrically integral and introduces no additional constants.

The smooth projective closure is an elliptic curve, with the displayed
point as origin.  Explicitly,
\[
 x=\frac{3\eta}{v+1},\qquad y=\frac{9(1-v)}{v+1}
 \quad\Longrightarrow\quad y^2=x^3-27.
\]
The normal Klein four subgroup of $A_4$ consists of translations by
the four two-torsion points of this elliptic curve.  The quotient
generator is
\[
 c:(v,\eta)\longmapsto(v,\zeta\eta).
\]
One can verify the translation assertion directly: translation by
$(3,0)$ sends $x$ to $3(x+6)/(x-3)$ and $y$ to
$-27y/(x-3)^2$; the resulting $v$ is $v_0$ in
\eqref{p2:eq:closure-roots}.  Conjugation by $c$ gives the other two
nontrivial translations.

\subsubsection{Why the deck action is the projective three-torsion action}
For the Legendre family with $\lambda=z^3$, put
\[
 \alpha(v)=\frac{(v^3+2)^2}{3(2v^3+1)}.
\]
Direct substitution gives
\[
 3\alpha^4-4(1+\lambda)\alpha^3+6\lambda\alpha^2-\lambda^2=0.
\]
Thus $\alpha$ is the $x$-coordinate of a nonzero point of order three,
or equivalently one of the four lines in the two-dimensional
$\mathbf F_3$-space of three-torsion.  The rational map $v\mapsto\alpha$
has degree six, whereas $v\mapsto z$ has degree four.  Therefore
$\alpha\notin\overline\Q_2(z)$: otherwise the degree six would be a
multiple of four.  The subgroup $C_3$ fixing $v$ is maximal in $A_4$,
so this already implies
\[
 \overline\Q_2(z,\alpha)=\overline\Q_2(v).
\]
The normal closures are equal.  Hence the $A_4$ cover just constructed
is precisely the splitting cover of the four projective three-torsion
lines, with their faithful permutation action.  Quadratic twists do
not change these lines and act trivially on $\operatorname{Sym}^4$.

\subsection{The semistable marked model of the true splitting cover}
Use the splitting cover and notation of the preceding subsection:
\[
\eta^3=4(v^3+1),\qquad e=\eta/\pi^2,\qquad
e^3=v^3+1,\qquad \pi^3=2,
\]
over $L=\Q_2(\zeta_3,\pi)$. Write $h=v^3$.
For completeness, the coordinate changes are explicit. Put
\[
 q=\frac{v^3+2}{v^3-1},\qquad
 \lambda=z^3=\frac{q^3(q+2)}{(q-1)(q+1)^3},\qquad
 \alpha=\frac{q^2}{q^2-1}
       =\frac{(v^3+2)^2}{3(2v^3+1)}.
\]
After quadratically twisting the Legendre curve by $\alpha$, the rational
three-torsion point and its tangent slope are
\[
 X_0=\alpha^2,\qquad
 Y_0=\frac{q^5}{(q-1)^3(q+1)^4},\qquad
 M_0=\frac{q(q^2+q+1)}{(q-1)(q+1)^2}.
\]
Translation by this flex gives $y^2+2M_0xy+2Y_0y=x^3$.
Scale first by $q/((q-1)(q+1))$ and then by $\pi$.
The resulting equation is
\begin{equation}\label{p2:eq:complete-tate}
y^2+\pi^2 A(v)xy+B(v)y=x^3,
\qquad
A=\frac{3(h^2+h+1)}{(h-1)(2h+1)},\quad
B=\frac{(h+2)^2}{(h-1)(2h+1)}.
\end{equation}
The quadratic twist disappears on $\Sym^4R^1$. Throughout this section,
scaling by $k$ means $x=k^2x'$, $y=k^3y'$, so that $(a_1,a_3)$ changes
to $(a_1/k,a_3/k^3)$.

The discriminant before pulling back to $e$ is
\[
\Delta=\frac{3^6h^2(h+1)^2(h+2)^6}{(h-1)^6(2h+1)^6}.
\]
On the $v$-line, the fourteen cusp fibres are $I_6$ at $v=0,\infty$ and
at the three roots of each of $h=-2$, $2h+1=0$, and $h=1$; the three
fibres at $h=-1$ are $I_2$. At either pole divisor $h=1$ or $2h+1=0$,
scaling by the inverse of its local parameter makes $a_1$ a unit,
$a_3$ vanish to order two, and the discriminant have order six.
Thus these poles do not conceal additive fibres.
The cubic cover $e^3=h+1$ ramifies at the three $I_2$ points, changing
them to $I_6$, and is unramified at all other cusps. Consequently all
36 cusp fibres on the true splitting cover are $I_6$.

The projective Fermat curve $C:e^3=v^3+1$ has good reduction at 2.
Its marked cusp sections form exactly nine clusters in the special fibre:
three over $v=0$, three at infinity, and three over $e=0$.
Use respectively $v$, $w=1/v$, and $e$ as local parameters. In each
cluster the four sections have parameter zero and three parameters whose
cubes are $\pm2$. The complementary coordinate is a unique power series
by Hensel's lemma, its derivative being three times a unit.
Blow up the nine disjoint ideals $(\text{parameter},\pi)$.
The resulting proper regular model $\mathcal B/\mathcal O_L$ has one
smooth genus-one main component and nine rational tails, all of
multiplicity one. Every tail has four disjoint cusp sections and meets
the main component at a fifth point. All 36 sections meet the smooth
locus of the special fibre and avoid the nine nodes.

This is the stable model of the curve with its full cusp divisor: each
rational tail has five special points. Uniqueness of the stable marked
model extends the deck group $G=A_4$ and the semilinear arithmetic action
$\sigma$ to $\mathcal B$. It also shows that the explicit blowups do not
depend on choices of labels within a cusp orbit.

\subsubsection{Good models over the nine nodes}
The following equations are obtained directly from
\eqref{p2:eq:complete-tate}. Their denominators are units in the indicated
node neighbourhoods.

At each node over zero put $vr=\pi$ and scale by $k=v^2$. Then
\[
a_1=r^2A(v),\qquad
a_3=\frac{(1+r^3)^2}{(v^3-1)(2v^3+1)}.
\]
At each infinity node put $w=1/v$, $wr=\pi$, and scale by $k=1/w$.
The coefficients become
\[
a_1=\frac{3r^2(1+w^3+w^6)}{(1-w^3)(1+r^3)},\qquad
a_3=\frac{(1+2w^3)^2}{(1-w^3)(1+r^3)}.
\]
At each node over $e=0$ put $er=\pi$, $h=e^3-1$, and scale by $k=1/e$.
Since $h-1=e^3-2=e^3(1-r^3)$, the coefficients are
\[
a_1=\frac{3r^2(h^2+h+1)}{(1-r^3)(2h+1)},\qquad
a_3=\frac{(h+2)^2}{(1-r^3)(2h+1)}.
\]
At the closed point of every node they reduce to $(a_1,a_3)=(0,1)$.
The Tate discriminant $a_3^3(a_1^3-27a_3)$ is therefore a unit.
These are smooth elliptic models over the node rings themselves, not
only good models over the generic points of the components.

\subsubsection{Integral lifts on all three kinds of tails}
Write $U=u^3$. At a zero tail put $v=\pi u$ and scale
\eqref{p2:eq:complete-tate} by $k=\pi^2$. One obtains
\[
a_1=\frac{3(4U^2+2U+1)}{(2U-1)(4U+1)},\qquad
a_3=\frac{(U+1)^2}{(2U-1)(4U+1)}.
\]
At an infinity tail put $w=\pi u$ and scale by
$k=1/(\pi(U+1))$. The integral lift is
\[
a_1=\frac{3(1+2U+4U^2)}{1-2U},\qquad
a_3=\frac{(1+4U)^2(U+1)^2}{1-2U}.
\]
At an $e=0$ tail put $e=\pi u$, so $h=2U-1$, and scale by
$k=1/(\pi(U-1))$. This gives
\[
a_1=\frac{3(h^2+h+1)}{2h+1},\qquad
a_3=\frac{(U-1)^2(h+2)^2}{2h+1}.
\]
Every displayed denominator is a unit along its finite special tail.
All three lifts reduce to the same ordinary elliptic family
\begin{equation}\label{p2:eq:common-tail}
E_u:\quad y^2+xy+(u^3+1)^2y=x^3,
\qquad \Delta=u^6(u^3+1)^6.
\end{equation}
It has four $I_6$ fibres, at $u=0$ and the three points $u^3=1$;
their $a_1$ is a unit. At the attaching point $u=\infty$, scaling by
$u^2$ gives good reduction, agreeing with the node charts.

On the main component, away from its nodes, scale
\eqref{p2:eq:complete-tate} by $v^2/e$. Its reduction is the constant
supersingular curve $E_0:y^2+y=x^3$. The node charts extend this good
reduction at all attaching points. The nine tails account for exactly
$9\cdot4=36$ fibres of type $I_6$, agreeing with the full generic cusp
divisor. Thus no horizontal collision or extra discriminant component
remains. The ratios of the main scale to the three node scales are
$1/e$, $1/(we)$, and $v^2$. The ratios of the three tail scales to
their node scales are $u^{-2}$, $u/(u^3+1)$, and $u/(u^3-1)$.
All are units on their respective overlaps; the latter two node charts
exclude the vanishing of their displayed denominators. These explicit
coordinate transitions glue the Weierstrass line bundle and cubic models.
At a horizontal cusp, $c_4$ is a unit and the exact discriminant divisor
is $\Delta=\text{unit}\cdot t^6$, with $t$ a local equation of its
section. The nodal deformation is therefore etale locally
$xy=\text{unit}\cdot t^6$; the unit can be absorbed in one coordinate.
The resulting
coefficient system has good reduction at every vertical node and
multiplicative reduction precisely along the 36 disjoint horizontal
sections. The functorial proper semistable/logarithmic cohomology
construction is consequently equivariant for $G$ and $\sigma$.

\subsection{From the charts to arithmetic tameness}

\begin{lemma}[Tameness from a semistable marked elliptic model]
\label{p2:lem:p2-log-bridge}
Let $K$ be a finite extension of $\mathbf Q_p$, let $L/K$ be a finite
 tamely ramified extension, and fix a prime $\ell\ne p$.  Let $U/K$ be
 a smooth geometrically connected curve, let $f:E\to U$ be an elliptic
 scheme, and put $\mathcal V=R^1f_*\mathbf Q_\ell$.
Suppose that there is a finite \'etale morphism $g:U'\to U_L$ of degree
 $d>0$ and an elliptic scheme $f':E'\to U'$ with
\[
 \operatorname{Sym}^4R^1f'_*\mathbf Q_\ell
   \simeq g^*\operatorname{Sym}^4\mathcal V.
\]
Suppose in addition that $U'$ has a proper semistable marked model
 $(\mathcal C,\mathcal D)/\mathcal O_L$ with the following properties.
\begin{enumerate}
\item The reduced horizontal divisor $\mathcal D$ is a disjoint union
 of \'etale sections in the smooth locus of $\mathcal C$, and
 $U'=\mathcal C_L\setminus\mathcal D_L$.
\item The elliptic scheme $E'$ extends to a proper flat family
 $\mathcal E\to\mathcal C$ which is smooth over
 $\mathcal C\setminus\mathcal D$.  In particular it is smooth at
 every node of the vertical fibre of $\mathcal C$.
\item Along $\mathcal D$ the family is semistable multiplicative:
 \'etale locally its singularities have equations
 $xy=t^m$ for some $m\ge1$, where $t=0$ defines a section of
 $\mathcal D$.  An invertible factor on the right is allowed.
\end{enumerate}
Then wild inertia acts trivially on
 $H^1(U_{\overline K},\operatorname{Sym}^4\mathcal V)$, and hence
 on its parabolic subspace
\[
 H^1_!(U_{\overline K},\operatorname{Sym}^4\mathcal V)
 :=\operatorname{im}\bigl(H^1_c(U_{\overline K},
 \operatorname{Sym}^4\mathcal V)\longrightarrow
 H^1(U_{\overline K},\operatorname{Sym}^4\mathcal V)\bigr).
\]
The same conclusion holds for every Galois-stable direct summand or
 subquotient of these representations.
\end{lemma}

\begin{proof}
Give $S=\operatorname{Spec}\mathcal O_L$ its standard divisorial log
 structure and $\mathcal C$ the log structure defined by its reduced
 vertical fibre together with $\mathcal D$.  This is a fine and
 saturated log-smooth curve over $S$.  Give $\mathcal E$ the log
 structure defined by the inverse image of this boundary, taking
 its fine and saturated toric model where necessary.  The stated
 node charts show that $\mathcal E\to\mathcal C$ is log smooth.
Here no condition $p\nmid m$ is needed for $xy=t^m$: the toric
 relation is $e_x+e_y=m e_t$, and the group cokernel relative to
 $\mathbf Z e_t$ is free, generated by $e_x=-e_y$.
An invertible factor can be absorbed in $x$ or $y$.  At a vertical
 base node the elliptic family is smooth, so its log chart is just
 the smooth pullback of $uv=\varpi$.  Thus composition gives a
 proper fine and saturated log-smooth model over $S$.

Form the fourfold fibre product in the category of fine and
 saturated log schemes,
\[
 \mathcal X=\mathcal E\mathbin{\times^{\mathrm{fs}}_{\mathcal C}}
 \mathcal E\mathbin{\times^{\mathrm{fs}}_{\mathcal C}}
 \mathcal E\mathbin{\times^{\mathrm{fs}}_{\mathcal C}}\mathcal E.
\]
Log smoothness is stable under this base change and composition.
The underlying scheme of $\mathcal X$ is proper over $S$: its
 toric saturation is finite over the ordinary proper fibre product.
The locus where the log structure is trivial is exactly
\[
 X^\circ=E'\times_{U'}E'\times_{U'}E'\times_{U'}E'.
\]
In particular it is an open subset of the generic fibre. The
 horizontal-boundary theorem \cite[Corollary~8.4.3, p.~310]{Illusie}
 applies to the trivial-log open of this proper log-smooth fine and
 saturated model. The original nearby-cycle and comparison statements are
 \cite[Variant~3.4.1, Proposition~4.2 and Remark~4.2.1]{Nakayama}.
Consequently wild inertia of $L$ is trivial on
 $H^q(X^\circ_{\overline L},\mathbf Z/\ell^n)$ for every $q,n$,
 and hence on its $\mathbf Q_\ell$ cohomology.  One must use the
 horizontal-boundary statement here; the generic log structure on
 the proper compactification is not trivial.

For completeness, the required coefficient cohomology is an
 explicit summand of ordinary degree-five cohomology of
 $X^\circ$.  Write $F:X^\circ\to U'$ and
 $\mathcal V'=R^1f'_*\mathbf Q_\ell$.  Let $\iota_i$ be inversion
 on factor $i$, and let $S_4$ permute the four factors.  The commuting
 idempotents
\[
 e_{\rm odd}=\prod_{i=1}^4\frac{1-\iota_i^*}{2},\qquad
 e_{\rm sym}=\frac1{4!}\sum_{\sigma\in S_4}
                 \operatorname{sgn}(\sigma)\sigma^*,\qquad
 e=e_{\rm odd}e_{\rm sym}
\]
are defined over $L$.  In relative K\"unneth cohomology,
 $e_{\rm odd}$ retains only the tensor of four degree-one factors.
Thus it kills $R^bF_*\mathbf Q_\ell$ for $b\ne4$ and has image
 $(\mathcal V')^{\otimes4}$ for $b=4$.  A transposition of two
 degree-one factors contributes the Koszul sign $-1$; the sign in
 $e_{\rm sym}$ therefore makes its action on this tensor power
 the ordinary symmetrizer.  It follows that
\[
 eR^bF_*\mathbf Q_\ell=
 \begin{cases}
 \operatorname{Sym}^4\mathcal V',&b=4,\\
 0,&b\ne4.
 \end{cases}
\]
Applying this projector to the Leray spectral sequence leaves
 only this one row.  Its canonical edge identification is
\[
 eH^5(X^\circ_{\overline L},\mathbf Q_\ell)
   \simeq H^1(U'_{\overline L},
                    \operatorname{Sym}^4\mathcal V').
\]
All maps are Galois equivariant, so the right-hand side is tame.

Finally the finite \'etale pullback and trace maps satisfy
\[
 \operatorname{Tr}_g\circ g^*=d\,\mathrm{id}
 \quad\text{on}\quad
 H^1(U_{\overline K},\operatorname{Sym}^4\mathcal V).
\]
After dividing the trace by $d$, pullback is a split injection
 into the tame coefficient cohomology just obtained.  The
 parabolic image is a Galois-stable subspace of ordinary
 cohomology.  Since $L/K$ is tamely ramified, its wild inertia
 subgroup equals the wild inertia subgroup of $K$ under the
 usual embedding $G_L\subset G_K$.  This proves all assertions.
\end{proof}


\begin{corollary}[Wild inertia on the cubic Legendre group]
\label{p2:cor:true-cover-swan}
For every $\ell\ne2$, the full parabolic group $\mathcal H_\ell$ and
both specified descents $W_{4,\ell}^{\pm}$ have zero Swan conductor at $2$.
\end{corollary}
\begin{proof}
Apply Lemma~\ref{p2:lem:p2-log-bridge} to the true splitting cover
\[
 z=\frac{v(v^3+2)}{2v^3+1},\qquad e^3=v^3+1,
 \qquad L=\Q_2(\zeta_3,\pi),\quad \pi^3=2.
\]
This constant extension has ramification index three and is tame.
If $h=v^3$, then
\[
 z^3-1=\frac{(h-1)^3(h+1)}{(2h+1)^3},\qquad
 \frac{dz}{dv}=\frac{2(v^3-1)^2}{(2v^3+1)^2}.
\]
The degree-four rational map ramifies only over the original cusps.
The Fermat projection ramifies at $v^3=-1$, which also maps into
$z^3=1$; it is unramified at the three points above infinity.
Removing all preimages of $\{0,\infty,\mu_3\}$ therefore gives a finite
etale degree-twelve cover. Its compactification has genus one and
its removed divisor has the 36 points already accounted for above.
The displayed integral charts and gluing maps verify each model
hypothesis of the lemma: separated etale horizontal sections, good
elliptic reduction at the nine vertical nodes, and horizontal
singularities $xy=\text{unit}\cdot t^6$. The twist parameter $\alpha=(h+2)^2/(3(2h+1))$ is invertible away
from these cusps, so it introduces no additional singular locus on
the generic open. The fourth symmetric power removes this quadratic
twist. The lemma proves tameness
of the full ordinary coefficient cohomology, hence of its parabolic
subspace and the two rational projected summands.
\end{proof}
\subsection{Geometric deck action and arithmetic inertia}
\subsubsection{The good main component and two different order-three actions}
Set
\[
 e=\eta/\pi^2,\qquad e^3=v^3+1.
\]
This Fermat model is smooth and proper over $\mathcal O_L$ and has
supersingular genus-one special fibre $\overline C$.  In particular
$\overline C[2](\overline{\mathbf F}_2)$ has just its origin.
The translations by two-torsion extend to this good model and
specialize to the identity.  Thus $K_4$ fixes the entire main
component pointwise.  In contrast, both $c$ and $\sigma$ act on its
base by $e\mapsto\zeta e$; for $\sigma$ this follows from
$\sigma(\eta)=\eta$ and $\sigma(\pi^2)=\zeta^2\pi^2$.

Their actions on the elliptic coefficient are different.  The exact
Tate equation derived above is
\[
 E^*: y_*^2+\pi^2 A(v)x_*y_*+B(v)y_*=x_*^3,
\quad
 A=\frac{3(h^2+h+1)}{(h-1)(2h+1)},\quad
 B=\frac{(h+2)^2}{(h-1)(2h+1)}.
\]
It was obtained by a scale $\pi$ after transformations defined over
$K_0(v)$, up to the harmless quadratic twist.  Hence
\[
 c:(x_*,y_*)\mapsto(x_*,y_*),\qquad
 \sigma:(x_*,y_*)\mapsto(\zeta x_*,y_*).
\]
On the main component use the further scale $k=v^2/e$, and write
$X=x_*/k^2$, $Y=y_*/k^3$.  Its special equation is the constant curve
\[
 E_0:Y^2+Y=X^3.
\]
For instance its $a_3$ before reduction is
$B(h+1)/h^2$, which reduces to one.  Transport of coordinates gives
the exact special-fibre actions
\begin{equation}\label{p2:eq:two-C3-actions}
 \begin{array}{c|cc}
 &\overline C&E_0\\\hline
 c&e\mapsto\zeta e&(X,Y)\mapsto(\zeta^2X,Y)\\
 \sigma&e\mapsto\zeta e&(X,Y)\mapsto(X,Y).
 \end{array}
\end{equation}
Possible signs from a choice of a square root in the preliminary
quadratic twist disappear on the fourth symmetric power.

The quotient of $\overline C$ by $e\mapsto\zeta e$ is the rational
$v$-line.  Since three is invertible in the $\ell$-adic coefficient
field, finite-map trace gives
\[
 H^1(\overline C,\Q_\ell)^{\sigma}=0.
\]
Thus, by \eqref{p2:eq:two-C3-actions}, the entire main component term
\[
 H^1(\overline C,\Q_\ell)\otimes
       \operatorname{Sym}^4H^1(E_0,\Q_\ell)
\]
has no $\sigma$-invariants.  This remains true for its $A_4$-invariant
part and for every subsequent projector.  The conclusion uses both
the base and coefficient actions, not the coefficient action alone.

\subsubsection{The projective quaternion action on the constant coefficient}
The full $A_4$ action on projective three-torsion remains faithful
after reduction: on a good elliptic model, three-torsion is finite
etale and its four projective lines remain distinct.  Since $K_4$
fixes the main base pointwise, its action on the constant elliptic
coefficient is therefore faithful in $\operatorname{Aut}(E_0)/\{\pm1\}$.
All these are zero-preserving elliptic isomorphisms.  Their projective
ambiguity is exactly the quadratic-twist sign already removed in
$\operatorname{Sym}^4$.

For completeness, the inverse image of this $K_4$ is the quaternion
group $Q_8$.  Its eight automorphisms are
\[
 (X,Y)\longmapsto(X+r,Y+r^2X+t),\qquad
 r\in\mathbf F_4,\quad t^2+t=r^3.
\]
For $r\ne0$ their squares are $[-1]:(X,Y)\mapsto(X,Y+1)$.
The order-three automorphism $(X,Y)\mapsto(\zeta^2X,Y)$ permutes
the three nonidentity elements of the projective $K_4$.  Together
these generate the binary tetrahedral group $Q_8\rtimes C_3$.

Let $M=\operatorname{Sym}^4 H^1(E_0,\Q_\ell)$.  The quaternion
character calculation gives
\begin{equation}\label{p2:eq:Q8-sym4}
 \dim M^{K_4}=\dim M^{Q_8}=2,\qquad
 M^{Q_8\rtimes C_3}=0.
\end{equation}
Indeed the traces on $\operatorname{Sym}^4$ are $5$ on the two
central elements and $1$ on the six elements of order four.  On
each of the sixteen remaining binary tetrahedral elements the
trace is $-1$.  Averaging gives $(10+6)/8=2$ and
$(10+6-16)/24=0$.  More precisely, if $c_0$ is either generator
of the quotient $C_3$, then
\[
 \operatorname{Tr}(c_0\mid M^{Q_8})
 =\frac18\sum_{q\in Q_8}\operatorname{Tr}(c_0q\mid M)=-1.
\]
Together with the absence of fixed vectors this shows that its
characteristic polynomial on $M^{K_4}$ is $X^2+X+1$, so its
eigenvalues are $\zeta,\zeta^2$.
This calculation is valid for every $\ell\ne2$, including $\ell=3$,
because coefficients have characteristic zero.

\subsubsection{Arithmetic action on the two node orbits}
The three attaching points over $z=0$ are $v=0$, $e^3=1$;
the three over $z=\infty$ are the three values of $e/v$ at infinity.
In each case $c$ and $\sigma$ permute these three points in the same
cycle.  The stabilizer in $A_4$ is $K_4$.  Use the good node scale
$k=v^2$ at zero and $k=1/w$ at infinity, where $w=1/v$.
These scales contain neither $e$ nor $\pi$.  Thus $c$ acts trivially
on the displayed elliptic node coordinates, whereas $\sigma$ acts
by $(X,Y)\mapsto(\zeta X,Y)$.

Identifying the $A_4$-invariant direct sum over either three-point
orbit with $M^{K_4}$ at one point, the arithmetic action is therefore
the coefficient action of $c^{-1}\sigma$.  It is the order-three
rotation $(X,Y)\mapsto(\zeta X,Y)$.  By \eqref{p2:eq:Q8-sym4}, neither
of these two two-dimensional node terms has arithmetic invariants.
Their duals and Tate twists have the same property, since the
$\ell$-adic Tate character is unramified at two.

At each of the remaining three main attaching points $e=0$, the
whole geometric $A_4$ is the stabilizer.  Its projective action is
the full tetrahedral action above, so its node coefficient invariant
space is zero by \eqref{p2:eq:Q8-sym4}.  The main degree-zero constant
coefficient invariant space is zero for the same reason.

Thus every main and node term surviving geometric descent has no
arithmetic eigenvalue one. The tail calculation and the exact
cohomological passage below complete the invariant calculation.
\subsection{The descended ordinary tails}
The degree-four tail quotient can be read from the exact rational map.
At zero, $v=\pi u$ gives
\[
\frac{z}{\pi^4}=\frac{u(u^3+1)}{4u^3+1}
\quad\longmapsto\quad Z=u^4+u.
\]
This is the translation quotient by $K_4=(\F_4,+)$. Since
$\lambda=z^3=16Z^3$, the reduced Legendre $j$-invariant is $1/Z^6$.
Infinity is its inversion counterpart.

At a tail over a cube root of unity, use $e=\pi u$ and $h=2u^3-1$.
The exact identity
\[
\lambda-1=\frac{(h-1)^3(h+1)}{(2h+1)^3}
\]
gives
\[
T=\frac{\lambda-1}{16}
=\frac{u^3(u^3-1)^3}{(4u^3-1)^3}
\quad\longmapsto\quad (u^4+u)^3.
\]
This is the degree-twelve quotient by the affine group
$A_4=(\F_4,+)\rtimes\F_4^\times$. Its reduced $j$-invariant is $1/T^2$.
The ordinary characteristic-two family
\[
\mathcal K_t:\quad y^2+xy=x^3+tx
\]
has $j=1/t^2$. Hence the two descended fourth-symmetric-power systems
are respectively its cubic pullback $t=Z^3$ and its ordinary version
$t=T$. Ordinary elliptic curves with the same $j$ differ by a quadratic
twist; this ambiguity disappears on $\Sym^4$.
Over the ordinary generic point, the zero-preserving automorphism group
is $\{\pm1\}$. Thus the identification is compatible with the descended
fourth-symmetric-power action, independently of the choice of elliptic
isomorphism.

These identifications include middle extensions at the missing points.
Indeed finite pushforward on the compact quotient curves is exact, and
finite-group invariants over $\Q_\ell$ are exact. The invariant pushforward
of a middle-extension coefficient system is the middle extension of its
generic descent, as may also be checked on the local inertia-invariant
stalks. This applies to the wildly ramified special translation quotients
as well: it uses characteristic-zero coefficients, not tameness of those
finite maps. Thus the component parabolic cohomology after deck descent
is exactly that of the displayed ordinary or cubic Kloosterman system.
The cited Kloosterman calculation gives zero for both of these parabolic
$\Sym^4$ groups, accounting for every descended tail middle term.
\subsection{Vanishing of the descended tail terms}

\begin{lemma}[The characteristic-two Kloosterman tails]
\label{p2:lem:p2-kl-tails}
Let $k=\overline{\mathbf F}_2$, let $\ell\ne2$, and let
 $f:E\to\mathbf G_{m,k}$ be the elliptic scheme
\[
 E_t:\quad y^2+xy=x^3+tx.
\]
Put $\mathcal F=\operatorname{Sym}^4R^1f_*\mathbf Q_\ell$, and
 let $[3]:\mathbf G_m\to\mathbf G_m$ send $Z$ to $Z^3$.
Then
\[
 H^1_!(\mathbf G_m,\mathcal F)=0,
 \qquad H^1_!(\mathbf G_m,[3]^*\mathcal F)=0.
\]
For the first sheaf the Swan conductor at infinity and the
 dimensions of the inertia-invariant spaces at $0,\infty$ are
 $(1,1,0)$; for the second they are $(3,1,2)$.
\end{lemma}

\begin{proof}
On the open set $x\ne0$, put $u=y/x$.  The elliptic equation becomes
 $u^2+u=x+t/x$, the Artin--Schreier cover defining the rank-two
 Kloosterman sheaf $\mathrm{Kl}_2$ in characteristic two.
The involution $u\mapsto u+1$ is elliptic inversion.  The two
 omitted points $x=0$ and the point at infinity are fixed by this
 involution, so their boundary cohomology is invariant.  Taking
 the odd part of the compact-support exact sequence identifies
 $R^1f_*\mathbf Q_\ell$ with $\mathrm{Kl}_2$.  This also fixes
 the trace convention: its Frobenius trace is minus the usual
 Kloosterman sum.  The elliptic realization is recorded in
 Haessig~\cite{Haessig}, proof of Theorem~2.2, citing Schoof--van der Vlugt,
 Proposition~3.1(4).

We use the characteristic-two local calculation of Yun~\cite{Yun},
 \S4.3.1 and Lemmas~4.3.3--4.3.4.  The image of inertia at
 infinity on $\mathrm{Kl}_2$ is the binary tetrahedral group
 $2.A_4$, with lower ramification filtration
\[
 D_0=2.A_4,\quad D_1=Q_8,\quad D_2=D_3=\{\pm1\},
 \quad D_4=1.
\]
In the fourth symmetric power this gives
 $\operatorname{Swan}_\infty(\mathcal F)=1$ and
 $\mathcal F^{I_\infty}=0$.  At zero, the local monodromy is a
 single unipotent block, so $\dim\mathcal F^{I_0}=1$.
The global geometric monodromy is $\mathrm{SL}_2$; therefore
 $H^0(\mathbf G_m,\mathcal F)=H^2_c(\mathbf G_m,\mathcal F)=0$.
The Grothendieck--Ogg--Shafarevich formula and the boundary
 exact sequence consequently give
\[
 \dim H^1_!(\mathbf G_m,\mathcal F)
 =\operatorname{Swan}_\infty(\mathcal F)
  -\dim\mathcal F^{I_0}-\dim\mathcal F^{I_\infty}=1-1-0=0.
\]

The cubic pullback kills the tame quotient $(2.A_4)/Q_8$ of
 the local image at infinity, leaving image $Q_8$.
On $\operatorname{Sym}^4$, its invariant plane is generated by
 $x^4+y^4$ and $x^2y^2$.  Tame pullback multiplies the Swan
 conductor by three.  Thus the corresponding local data are
 $\operatorname{Swan}_\infty=3$, invariant dimension $1$ at zero,
 and invariant dimension $2$ at infinity.  A finite-index
 restriction still has geometric Zariski closure $\mathrm{SL}_2$,
 so the global $H^0$ and $H^2_c$ again vanish.  The same formula
 now reads $3-1-2=0$.
\end{proof}

\begin{remark}
The lemma applies to the ordinary tail $t=(\lambda-1)/16$ and
 to the cubic tails $\lambda=16Z^3$ supplied by the explicit
 Legendre coordinate changes, after quadratic twists are removed
 by $\operatorname{Sym}^4$.  The statement concerns parabolic
 cohomology of these characteristic-two families.  Identifying
 their contributions to the original dyadic inertia invariants
 still requires the component and node maps of the semistable
 model; the vanishing alone is not a local-factor calculation.
\end{remark}
\subsection{Equivariant nearby cycles and parabolic descent}

We apply the preceding calculations to the proper marked semistable
model $\mathcal B/\mathcal O_L$ of the actual splitting cover
$\eta^3=4(v^3+1)$. Its main component is $C:e^3=v^3+1$, where
$e=\eta/\pi^2$. Write $G=A_4$ for the geometric deck group,
$c:e\mapsto\zeta e$ for the indicated geometric order-three operation,
and $\sigma(\pi)=\zeta\pi$ for the arithmetic descent operation.
The two operations on the coefficient are those computed in the
arithmetic-action calculation above; they must not be identified.

The marked model has a genus-one main component and nine rational tails,
each with four marked cusp sections. It is a stable marked curve.
Consequently the deck and arithmetic descent operations extend to it
by uniqueness of the stable marked model. Let $\mathcal D$ be the union
of the horizontal cusp sections, and write
\[
 j:\mathcal B\setminus\mathcal D\hookrightarrow\mathcal B,
 \qquad
 \mathcal F=\operatorname{Sym}^4R^1f_*\mathbf Q_\ell,
 \qquad \mathcal J=j_*\mathcal F.
\]
Here $f$ is the good elliptic family on the indicated open, supplied by
the integral charts. We assert equivariance for $\mathcal F$, not for
an arbitrary choice of the preliminary quadratic twist of the elliptic
family. On the generic open, $\mathcal F$ is canonically the pullback
of the original fourth symmetric-power coefficient. Its deck and
arithmetic identifications extend to the lisse sheaf on
$\mathcal B\setminus\mathcal D$: restriction of lisse sheaves to a
dense open of a normal scheme is fully faithful. They then extend to
$\mathcal J$ by functoriality of $j_*$. In particular the following
sequences carry the actual $G$ and arithmetic actions.

\subsubsection{The local nearby-cycle calculation}

Let $\mathcal B_a$ denote the normalized special-fibre components, and
let $\mathcal J_a$ be their coefficient middle extensions across their
horizontal cusps. For each of the nine vertical nodes $x$, put
$H_x=\mathcal F_x$ and set
\[
 \mathcal N=\bigoplus_{x\text{ a vertical node}}H_x.
\]
The good elliptic charts at the nodes are essential here. The
semistable-curve calculation gives
\begin{equation}\label{p2:eq:complete-p2-nearby}
 R^0\Psi(\mathcal J_{\bar\eta})=\mathcal J_{\bar s},\qquad
 R^1\Psi(\mathcal J_{\bar\eta})
     =\bigoplus_x(H_x(-1))_x,\qquad
 R^q\Psi(\mathcal J_{\bar\eta})=0\quad(q>1).
\end{equation}
At a vertical node the coefficient is lisse on the full ring
$uv=\pi$. Thus the assertion is the constant-sheaf node calculation
tensored with $H_x$; see \cite[Corollary~8.4.4]{Illusie} for that
semistable specialization of the nearby-cycle formula.

There are no extra vertical vanishing cycles at the horizontal cusps.
The base is smooth there, the marked section is a relative smooth
divisor, and the elliptic family is multiplicative with local parameter
$q=t^6$ times a unit. After strict henselization the branches split and
the unit has every prime-to-$2$ root. The coefficient therefore has its
usual tame unipotent monodromy in $t$, and its horizontal invariant
sheaf $\mathcal F^{I_t}$ is lisse over the trait.
More precisely, \cite[Expose~XIII, Lemmas~2.1.10--2.1.11]{SGA7}
gives base change for $Rj_*\mathcal F$ and
$R\Phi(j_!\mathcal F)=0$ for a lisse sheaf tamely ramified along a
relative normal-crossings divisor. Apply also smooth local acyclicity
to the invariant sheaf on the section. The exact sequence
\[
 0\longrightarrow j_!\mathcal F\longrightarrow j_*\mathcal F
 \longrightarrow i_*\mathcal F^{I_t}\longrightarrow0
\]
then gives $R\Phi(j_*\mathcal F)=0$ at that section. The statements
apply at finite prime-to-$2$ coefficient levels and then to
$\mathbf Q_\ell$ coefficients. Away from the cusps and vertical nodes,
ordinary smooth local acyclicity applies. These observations prove
\eqref{p2:eq:complete-p2-nearby} and identify the restrictions of
$\mathcal J_{\bar s}$ with the component middle extensions.

Proper base change and the five-term nearby-cycle sequence yield
\begin{equation}\label{p2:eq:complete-p2-five-term}
 0\longrightarrow H^1(\mathcal B_{\bar s},\mathcal J_{\bar s})
 \longrightarrow H^1(\mathcal B_{\bar\eta},\mathcal J_{\bar\eta})
 \longrightarrow\mathcal N(-1)
 \longrightarrow H^2(\mathcal B_{\bar s},\mathcal J_{\bar s}).
\end{equation}
For the proper comparison, see
\cite[Expose~XIII, \S2.1.8, especially (2.1.8.9)]{SGA7}.
We only use that the quotient of the second term by the first embeds
in $\mathcal N(-1)$; surjectivity onto the node module is unnecessary.

For normalization $\nu:\coprod_a\mathcal B_a\to\mathcal B_{\bar s}$,
lissity at the nodes gives
\[
 0\longrightarrow\mathcal J_{\bar s}
 \longrightarrow\nu_*\bigoplus_a\mathcal J_a
 \longrightarrow\bigoplus_x(H_x)_x\longrightarrow0.
\]
The last map is the difference of the two restrictions. Orient every
edge from the unique genus-one main component toward its tail. Both
$G$ and $\sigma$ preserve this orientation, so there is no additional
sign in their action on the node module. The cohomology sequence gives
\begin{equation}\label{p2:eq:complete-p2-normalization}
 \bigoplus_a H^0(\mathcal B_a,\mathcal J_a)
 \longrightarrow\mathcal N
 \longrightarrow H^1(\mathcal B_{\bar s},\mathcal J_{\bar s})
 \longrightarrow\bigoplus_a H^1(\mathcal B_a,\mathcal J_a)
 \longrightarrow0.
\end{equation}
Both sequences are equivariant. Taking $G$-invariants is exact over
$\mathbf Q_\ell$, also when $\ell=3$.

\subsubsection{Finite descent, including the wild special-tail quotients}

We recall why taking the indicated stabilizer invariants on a special
tail computes its descended middle cohomology even when the quotient
is wildly ramified at the attaching point. Let
$q:Y\to X=Y/H$ be any of the finite special-tail quotient maps.
On the dense open where the coefficient descends, let $F_X$ be the
descended sheaf and let $j_Y,j_X$ be the respective open immersions.
Finite pushforward and exactness of rational $H$-invariants give
\[
 \bigl(q_*j_{Y*}q^*F_X\bigr)^H=j_{X*}F_X.
\]
This can be checked directly on sections: extension by $j_*$ and
finite pushforward commute, and invariant sections on the finite
cover are precisely sections downstairs. Hence
\[
 H^1(Y,j_{Y*}q^*F_X)^H=H^1(X,j_{X*}F_X).
\]
No tameness of $q$ at the omitted point is required. In the present
case the maps on the open tails are the displayed separable quotients
$Z=u^4+u$ and $T=(u^4+u)^3$. Their ramification at infinity is thus
handled by the middle extension, not discarded. The same argument
justifies generic finite Galois descent of parabolic cohomology.

The preceding component and group calculations now supply exactly the
following inputs:
\begin{enumerate}
\item The main-component cohomology has no $\sigma$-eigenvalue $1$:
$\sigma$ acts on its Fermat-base $H^1$ with the two nontrivial cubic
characters and acts trivially on the constant elliptic coefficient.
\item All descended tail middle-cohomology groups vanish, by their
identification with the ordinary and cubic-pullback Kloosterman
families and the stated local-monodromy calculation.
\item $\mathcal N^G$ consists of the two two-dimensional node spaces
over $0$ and $\infty$. Each has arithmetic eigenvalues
$\zeta,\zeta^2$; the three other node spaces vanish after $G$-invariants.
The Tate twist $(-1)$ is unramified at $2$.
\end{enumerate}
Consequently neither $\mathcal N^G$ nor
$(\bigoplus_a H^1(\mathcal B_a,\mathcal J_a))^G$ has eigenvalue $1$.
Equation~\eqref{p2:eq:complete-p2-normalization} gives the same conclusion
for $H^1(\mathcal B_{\bar s},\mathcal J_{\bar s})^G$.
Equation~\eqref{p2:eq:complete-p2-five-term} then gives it for
$H^1(\mathcal B_{\bar\eta},\mathcal J_{\bar\eta})^G$.
Absence of an eigenvalue passes to subspaces, quotients and extensions,
so no degeneration assertion or vanishing of graph monodromy is used.

Finally, $H^1$ of the coefficient middle extension on a proper smooth
curve is canonically its parabolic cohomology. Generic finite Galois
descent therefore identifies the last $G$-invariant group with the full
original cubic-Legendre parabolic representation. An inertia-invariant
vector would be fixed by the chosen arithmetic inertia lift $\sigma$,
which has no eigenvalue $1$ here. Both specified rank-four summands
thus have zero inertia invariants. Combining this with Corollary~\ref{p2:cor:true-cover-swan} gives, for every $\ell\ne2$,
\[
 d_{2,\ell}^+=d_{2,\ell}^-=0,\qquad
 P^+_{2,\ell}(T)=P^-_{2,\ell}(T)=1,\qquad
 a_2(W^+_{4,\ell})=a_2(W^-_{4,\ell})=4.
\]


This proves Theorem~\ref{p2:thm:dyadic-local-data}. In particular the
arithmetic constant in the splitting cover is retained throughout;
replacing $\eta^3=4(v^3+1)$ by a geometrically isomorphic equation
without tracking that constant would not determine this inertia action.

\section{Conductor products and coefficient-prime scope}\label{sec:conductors}
For a realization $V_\ell$ and a prime $p\ne\ell$, use
\[
 P_p(V_\ell,T)=\det(1-\Frob_pT\mid V_\ell^{I_p}),\qquad
 a_p(V_\ell)=\dim V_\ell-\dim V_\ell^{I_p}+\Sw_p(V_\ell).
\]
The local proofs are carried out directly with arbitrary coefficient prime
different from the residue characteristic. They give
\[
\begin{array}{c|ccc|ccc}
 &d_2&\Sw_2&a_2&d_3&\Sw_3&a_3\\\hline
 W_{4,\ell}^+&0&0&4&1&0&3\\
 W_{4,\ell}^-&0&0&4&0&0&4
\end{array}
\]
with the $p$-column interpreted only for $\ell\ne p$. At $3$ the invariant
line of the plus descent has geometric Frobenius eigenvalue $3$.
Thus $P_2^+=P_2^-=P_3^-=1$ and $P_3^+=1-3T$.

For $p\ge5$ and $\ell\ne p$, good-prime unramifiedness follows by the
relative tame-specialization argument of Section~\ref{sec:geometry}.
To justify its coefficient change explicitly, use the integral sheaf on
the full cubic cover before applying rational projectors. At boundary points of both generic and special geometric fibres,
the Tate-module description gives prime-to-$p$ unipotent local
monodromy; at infinity the additional quadratic twist is also tame.
The cubic Kummer cover is tame since $p\ne3$. Thus every finite
$\ell$-power coefficient level is relatively tame. The same specialization,
duality and parabolic-image comparison then applies, and rational
projectors pass to its summands. This uses the local Tate descriptions,
not a claim that the order of $\GL_2(\Z/\ell^n)$ is prime to $p$.
Hence, for every $\ell\ne2,3$,
\[
 N^{(\ell)}(W_{4,\ell}^+)=\prod_{p\ne\ell}p^{a_p}=2^4\,3^3=432,
 \qquad
 N^{(\ell)}(W_{4,\ell}^-)=2^4\,3^4=1296.
\]
Equivalently these are the conductor products obtained by choosing an
auxiliary coefficient prime away from each residue characteristic.
Independence of the displayed bad-prime factors and exponents follows
from the explicit calculations, without invoking a general bad-prime
compatibility theorem. No assertion about the usual inertia invariants
of the $p$-adic realization at $p=\ell$ is being substituted for an
away-from-$\ell$ Euler factor. Nor does this calculation identify the
entire Weil--Deligne representation in all coefficient characteristics.

The numerical conductors formerly proposed for the two recipes therefore
agree with the conductor products of the specified descents. This does
not establish the functional-equation sign, analytic continuation, or
the accuracy of the historical central-value computation.

\section{Strong irreducibility from local monodromy}\label{sec:irreducibility}
\paragraph{Conventions and inputs.}
Fix a coefficient prime $\ell\ne3$. Write $G_\Q=\Gal(\overline\Q/\Q)$ and
$\tau_\ell$ for the character of $\Q_\ell(-1)$, so that
$\tau_\ell(\Frob_p)=p$ for geometric Frobenius at $p\ne\ell$.
To avoid opposite Hodge--Tate sign conventions, call $h$ a \emph{Tate
exponent} when a Hodge--Tate summand is $\C_\ell(-h)$.
Thus $\Q_\ell(-1)$ has Tate exponent $1$.

The inputs for $V=W_{4,\ell}^+$ are:
\begin{enumerate}
\item $V$ has dimension four over $\Q_\ell$, is unramified outside a finite
set, and is de Rham at $\ell$.
\item Its Tate exponents are $\{0,0,5,5\}$. These follow from the specified
smooth projective Scholl realization of
Theorem~\ref{thm:scholl-realization}, its rational projector and the Hodge
numbers $h^{5,0}=h^{0,5}=2$, using de Rham and Hodge--Tate comparison.
\item At $3$, inertia is unipotent with a monodromy logarithm $N$ having
one Jordan block of length four. With $F$ geometric Frobenius,
\[
 FNF^{-1}=3^{-1}N,
\]
and the eigenvalues on the successive quotients of
$0\subset\ker N\subset\ker N^2\subset\ker N^3\subset\ker N^4=V$
are $3,9,27,81$, respectively.
\end{enumerate}
The third input is the full arithmetic conclusion of
Section~\ref{sec:local-three}, not an inference from the conductor alone. The same construction gives
$W_{4,\ell}^-=W_{4,\ell}^+\otimes\chi_{-3}$.

\begin{lemma}[Rank-one characters over $\Q$]\label{irr:rankone}
Let $\delta:G_\Q\to\Q_\ell^\times$ be a continuous character, unramified
outside a finite set, and Hodge--Tate at $\ell$ with Tate exponent $m$.
Then $\delta=\epsilon\tau_\ell^m$ for a finite-order character $\epsilon$.
\end{lemma}
\begin{proof}
Tate's theorem identifies a Hodge--Tate abelian local representation with
a locally algebraic one when its inertia action is semisimple. A character
automatically satisfies the latter condition. The precise primary reference
is Serre \cite[Chapter~III, \S1.2, Theorem (Tate), printed p.~III-7;
proof in the Appendix to Chapter~III]{SerreRankOne}; the definition in \S1.1
identifies the local algebraic power. In the present $\Q_\ell$-valued
rank-one case, it follows that $\delta\tau_\ell^{-m}$ has finite image
on inertia at $\ell$.

Put $\epsilon=\delta\tau_\ell^{-m}$. By Kronecker--Weber, an abelian
character of $G_\Q$ is a continuous character of
$\widehat\Z^\times=\prod_p\Z_p^\times$. Outside a finite set its
restrictions to the unit factors are trivial, by the unramified hypothesis.
For each remaining $p\ne\ell$, its restriction to $\Z_p^\times$ has
finite image: its open pro-$p$ subgroup maps into $\Q_\ell^\times$,
whose compact subgroups have an open pro-$\ell$ subgroup, so that image
is finite. At $\ell$ finiteness was just proved. Thus the entire image is
contained in the product of finitely many finite groups, and is finite.
The exponent $m$ is the stated Tate exponent, since finite characters
have Tate exponent zero.
\end{proof}

\begin{lemma}[The single-block flag]\label{irr:flag}
Let $k$ be a field of characteristic zero and let $N$ on $k^4$ have one
nilpotent Jordan block. Its invariant subspaces are exactly
$\ker N^r$ for $0\le r\le4$.
\end{lemma}
\begin{proof}
As a $k[t]$-module, with $t$ acting as $N$, this space is
$k[t]/(t^4)$. Its submodules are the ideals $(t^j)/(t^4)$,
equivalently the displayed kernels. This proof remains valid after
arbitrary scalar extension.
\end{proof}

\begin{theorem}\label{irr:absolute}
Under the stated inputs, both $W_{4,\ell}^+$ and $W_{4,\ell}^-$ are
absolutely irreducible for every $\ell\ne3$.
\end{theorem}
\begin{proof}
It suffices to treat $V=W_{4,\ell}^+$. Suppose
$0\ne U\subsetneq V\otimes_{\Q_\ell}\overline\Q_\ell$ is a
$G_\Q$-stable subspace, of dimension $r\in\{1,2,3\}$.
Its finitely many basis coordinates belong to a finite coefficient
extension, so there is no issue of defining a continuous representation
over an infinite coefficient field. There is in fact a stronger descent
argument in this case.

Some inertia element at $3$ acts as $\exp(cN)$ with $c\ne0$.
Since its logarithm is a finite polynomial in $\exp(cN)-1$, stability
under that element makes $U$ stable under $N$. Lemma~\ref{irr:flag}
gives
\[
 U=(\ker N^r)\otimes_{\Q_\ell}\overline\Q_\ell.
\]
Set $U_0=\ker N^r\subset V$. This is defined over $\Q_\ell$ and is
$G_\Q$-stable: extension of scalars makes every $\rho(g)$ preserve it,
and the original matrices are $\Q_\ell$-linear. Hence its determinant
$\delta=\det(U_0)$ is a $\Q_\ell$-valued character.

The category of de Rham representations is closed under subobjects,
quotients and exterior powers; its comparison functor is exact and
strict for the filtrations. Thus $U_0$ and $\delta$ are de Rham, the Tate
exponents of $U_0$ form a submultiset of $\{0,0,5,5\}$, and the
Tate exponent $m$ of $\delta$ is their sum. In particular
\[
 m\in5\Z.
\]
For clarity, the dimension part of closure under subobjects follows already
from left exactness of $D_{\mathrm{dR}}$ and
$\dim D_{\mathrm{dR}}(M)\le\dim M$: apply this to $U_0$, $V$ and
$V/U_0$, with equality for $V$. The usual strictness of de Rham
comparison supplies the asserted filtered submultiset statement.

The character $\delta$ is unramified outside a finite set, so
Lemma~\ref{irr:rankone} gives $\delta=\epsilon\tau_\ell^m$.
Since $U_0$ has unipotent inertia at $3$, its determinant is unramified
there, as is $\epsilon$. On the other hand the explicitly ordered
monodromy flag gives
\[
 \delta(\Frob_3)=3^{1+2+\cdots+r}=3^{r(r+1)/2}.
\]
Therefore
\[
 \epsilon(\Frob_3)=3^{r(r+1)/2-m}.
\]
The left side is a root of unity. A rational integer power of $3$
is a root of unity only when its exponent is zero. Consequently
$m=r(r+1)/2$. For $r=1,2,3$ these values are $1,3,6$, respectively,
none of which belongs to $5\Z$. This contradiction proves absolute
irreducibility of $V$. Tensoring with the one-dimensional character
$\chi_{-3}$ preserves absolute irreducibility and proves the minus case.
\end{proof}

\begin{corollary}[Restriction to finite extensions]\label{irr:strong}
For every finite extension $L/\Q$, the restrictions of either
$W_{4,\ell}^\pm$ to $G_L$ are absolutely irreducible, for $\ell\ne3$.
\end{corollary}
\begin{proof}
It suffices again to treat the plus representation. If a proper nonzero
subspace $U$ is stable under $G_L$, replace $L$ by a finite Galois
extension containing it and put $H=G_L\triangleleft G_\Q$.
The intersection $H\cap I_3$ has finite index in $I_3$, hence its image
under the nonzero tame monodromy parameter still contains a nonzero
element. Consequently every $H$-stable subspace is $N$-stable.
Normality makes every $gU$ also $H$-stable, and all have dimension $r$.
By Lemma~\ref{irr:flag}, all equal $\ker N^r$ after scalar extension.
Thus $U$ is $G_\Q$-stable, contradicting Theorem~\ref{irr:absolute}.
\end{proof}

\paragraph{Limits and consequences.}
The argument uses the written Scholl/Hodge and local-at-$3$ results as
inputs; it does not re-prove their constructions. It makes no assertion
for the coefficient prime $\ell=3$, since the local monodromy calculation
above uses $\ell\ne3$. No conjecture of semisimplicity is used.
Proper induction from a smaller-dimensional representation over a finite
extension is excluded by Corollary~\ref{irr:strong}. An ordinary tensor
product of two rank-two continuous $\ell$-adic representations over finite
coefficient extensions is also excluded:
after a finite local extension its monodromy operator would have form
$N_1\otimes1+1\otimes N_2$, with $N_i^2=0$, hence its cube would be
zero, whereas the actual $N^3$ is nonzero. These exclusions are independent
of the dyadic calculation in Section~\ref{sec:local-two} and do not
establish automorphy or determine the Sato--Tate group.


\section{Current claims and their dependencies}\label{sec:claims}
``Written proof'' below describes the arguments in this manuscript.
It does not mean external peer review or formal verification. The
coefficient-prime restrictions are part of the conclusions.

\small
\begin{longtable}{>{\raggedright\arraybackslash}p{0.23\textwidth}>{\raggedright\arraybackslash}p{0.35\textwidth}>{\raggedright\arraybackslash}p{0.32\textwidth}}
\toprule
Conclusion & Proof location and main inputs & Scope or remaining limit\\
\midrule
\endfirsthead
\toprule Conclusion & Proof location and main inputs & Scope or remaining limit\\
\midrule\endhead
Ranks and rational descents & Sections~\ref{sec:geometry}, \ref{sec:descent}:
parabolic dimension, literal inversion and deck projectors &
Specified arithmetic normalization; $W_4^-=W_4^+\otimes\chi_{-3}$.\\[4pt]
Pairing and good Euler factors & Sections~\ref{sec:descent}, \ref{sec:inert}:
cup product, exact finite sums and inert boundary term &
Exact formulas at $p\ge5$; no certification of every stored numerical row.\\[4pt]
Good-prime weight five & Corollary~\ref{cor:descent-good-prime-purity}:
Deligne's weight theorem and specialization &
Independent of empirical reciprocal-quartic tests.\\[4pt]
Rational Chow realization and Hodge numbers &
Theorem~\ref{thm:scholl-realization}: full-level cover, finite projectors,
Scholl comparison & All $\ell$; $h^{5,0}=h^{0,5}=2$; no other Hodge types.\\[4pt]
Local data at $3$ & Section~\ref{sec:local-three}: quotient family,
collision cohomology, ordered Frobenius flag &
Every $\ell\ne3$; plus $P_3=1-3T$, $a_3=3$; minus $P_3=1$, $a_3=4$.\\[4pt]
Local data at $2$ & Section~\ref{sec:local-two}: exact Galois closure,
semistable charts, tame comparison, tails and nodes &
Every $\ell\ne2$; both $d_2=0$, $P_2=1$, $\Sw_2=0$, $a_2=4$.\\[4pt]
Conductor products & Section~\ref{sec:conductors}: both local
calculations and good-prime unramifiedness &
Plus $432$, minus $1296$; for $\ell\ne2,3$, these are the prime-to-$\ell$
Artin conductors.\\[4pt]
Strong absolute irreducibility & Section~\ref{sec:irreducibility}:
single-block monodromy at $3$, Hodge--Tate exponents and rank-one
characters over $\Q$ & Every $\ell\ne3$, after every finite
number-field extension; coefficient prime $3$ is not covered.\\[4pt]
Cubic second-moment comparison & Section~\ref{sec:cubic}:
geometric trace/determinant, stated comparison theorem, seven exact rows
and Sturm bound $24$ & Separate rank-two object; bounded exact
certificate retained; orders $4,6$ do not inherit this proof.\\
\bottomrule
\end{longtable}
\normalsize

The local factor and exponent independence asserted here is obtained
directly for all allowable auxiliary primes. A full comparison of the
bad-prime Weil--Deligne representations, including a treatment at equal
coefficient and residue characteristic, is a different statement.
The remaining research questions are listed in Section~\ref{sec:limits}.

\paragraph{Sign dictionary.}
At an inert good prime, the literal plus descent $I=+1$ has trace
$-S_p-1$ and conductor $432$. The literal minus descent $I=-1$ has trace
$S_p+1$ and conductor $1296$. Thus an older numerical label ``plus sign''
for the finite-sum formula means the latter object; it is not the
literal plus eigenspace of $I$.

\paragraph{Review priorities.}
The most useful independent review targets are the arithmetic constant
in the dyadic closure, the coefficient action on the main component and
nodes, the horizontal-boundary nearby-cycle hypotheses, and the
local-at-$3$ global-to-local map with its Frobenius normalization.
The strong-irreducibility deduction depends on the complete local flag,
not merely on its conductor exponent. These are review targets within
the stated proof, not newly claimed external endorsements.

\section{Numerical evidence, corrections and open questions}\label{sec:limits}
Theorem~\ref{thm:main} identifies the \emph{exact good-prime recipe}.
It does not turn each stored approximation to that recipe into exact data.
The code supplement preserves the earlier correction package and its logs
unchanged. No archived research program, prime scan, certificate generator
or numerical $L$-function evaluator was executed for this consolidation.
The previously supplied standalone implementation rechecked only the
bounded cubic certificate and Sturm comparison in Section~\ref{sec:cubic};
its results are retained separately in the supplied supplement. The local continuations used only the small exact algebra checks recorded
in their supplements. No archived point-counting scan or numerical
$L$-function experiment has been rerun for this consolidation.

\subsection{Evidence retained from the earlier computations}
The carried record contains direct exact small-prime counts, an exact
finite modular-comparison certificate, rounded FFT-derived larger tables,
and high-precision evaluations using those tables. These are different
evidence classes. An exact consistency check on rounded input does not
certify how the input was obtained. The proof of the reciprocal template
also does not validate every historical value of its first two coefficients.

The historical $1296$ model imposed $P_2(T)=P_3(T)=1$.
The later $432$ twist in the retained N14/N15 record used
$P_2(T)=1$ and $P_3(T)=1-3T$. The blanket statement in version 2
that both experiments imposed $P_2=P_3=1$ was therefore too broad.
The new local proofs establish these factors and conductor products for
$W_4^-$ and $W_4^+$, respectively; they do not use the numerical fits.
The historical phrase ``plus-sign recipe'' refers to inert trace $S_p+1$,
which is the \emph{literal minus descent} $W_4^-$ of this manuscript.
For that $1296$ model the recorded functional-equation sign is $+1$
and the central value is approximately $2.2413552$. Those two statements
remain numerical observations, not consequences of the local proofs.
The candidate completion is
\[
 \Lambda(s)=N^{s/2}\Gamma_{\C}(s)^2\mathcal L(s),\qquad
 \Gamma_{\C}(s)=2(2\pi)^{-s}\Gamma(s).
\]
The associated theta-kernel Mellin transform used in the code has the
normalization $\widehat\Lambda=\Lambda/4$. Working precision and small
finite-sum defects do not supply an interval enclosure or an infinite-tail
bound. The historical gamma-factor shape agrees with the independent Hodge
calculation in Section~\ref{sec:scholl}. The conductor is now independently established by the local arguments.
The gamma-factor agreement and conductor calculation do not prove the
functional-equation sign or analytic continuation.

Likewise, a symmetrized evaluator
\[
 \sum_{n\le N}a_n\{G(s,n,t_0)+\epsilon G(6-s,n,1/t_0)\}
\]
has built-in symmetry at $t_0=1$. For $s=3$ and $\epsilon=-1$ it
vanishes for any coefficient sequence. This is not an independent
analytic-rank test. The historical N18 failure caused by comparing an
approximately fifty-digit reference against a seventy-digit requirement
is retained as a failure; other reported agreements are not promoted to
certified analytic digits.

\subsection{Two exact algebraic checks and their limits}
\begin{lemma}[Reciprocal-quartic modulus criterion]\label{lem:weil}
Let $P>0$, $c_1,c_2\in\R$,
$D=c_1^2-4(c_2-2P)$ and $B=16P-c_1^2-D$.
Every root of $X^4-c_1X^3+c_2X^2-Pc_1X+P^2$ has modulus
$\sqrt P$ if and only if
\[
 D\ge0,\qquad B\ge0,\qquad B^2\ge4c_1^2D.
\]
\end{lemma}
\begin{proof}
Put $y=X+P/X$. The two $y$-values are the roots of
$y^2-c_1y+c_2-2P$. The required condition is that they are real
and both lie in $[-2\sqrt P,2\sqrt P]$, equivalently
$D\ge0$ and $|c_1|+\sqrt D\le4\sqrt P$.
Squaring with the nonnegative signs retained gives the stated test.
\end{proof}
This repairs the earlier floor-square-root validator; it tests specified
coefficients. Good-prime purity of the exact geometric factors follows
independently from Corollary~\ref{cor:descent-good-prime-purity};
it does not certify any stored coefficient table.

If the same quartic factors as
$(X^2-aX+P)(X^2-bX+P)$, coefficient comparison gives
$D=(a-b)^2$. A nonsquare rational $D$ excludes this model with
$a,b\in\Q$. If $a,b$ are conjugate in a fixed quadratic field,
every nonzero $D$ has that field's square class. Neither this limited
obstruction nor empirical trace distributions establishes general
irreducibility, rules out every inner twist, or determines the Sato--Tate group.
The independent proof in Section~\ref{sec:irreducibility} now establishes
strong absolute irreducibility for $\ell\ne3$; it does not rely on this test.

\subsection{Remaining mathematical questions}
The bad-prime factors and conductor exponents have been determined in
Sections~\ref{sec:local-three}--\ref{sec:conductors}. Their auxiliary-prime
scope is recorded there; they are no longer open obligations of this draft.
The following questions remain separate.
\begin{enumerate}
\item The strong-irreducibility argument uses $\ell\ne3$. Its extension
at coefficient prime $3$, a full bad-prime Weil--Deligne comparison, and
precise Sato--Tate or inner-twist classifications are not established here.
\item Identify any automorphic representation with all required
hypotheses checked. The manuscript does not establish quaternion
multiplication or transfer an adjacent automorphy theorem.
\item Establish analytic continuation and the functional equation by a
theorem, including its sign. A root-number calculation requires its own
local epsilon-factor inputs; the conductor alone does not determine it.
\item Fix rational period structures before interpreting critical
values or determinant ratios arithmetically. The retained numerical
central values are not certified analytic evaluations.
\item For character orders $4,6$ in the older second-moment work,
supply their separate geometric inputs. Their finite comparison rows do
not follow from the cubic-only proof as all-prime theorems.
\end{enumerate}
These are limits of the present conclusions, not instructions to pursue
all of them before this manuscript can be independently reviewed.
The rank-four object is not the rank-two $H^1$ of an elliptic curve, even
after a Tate twist. No comparison controlling the Mordell--Weil rank or
BSD leading coefficient of a specified elliptic curve has been established
here. A future motivic special-value formulation would require its own
precise twist, arithmetic groups and comparison theorem.


\appendix
\section{Known quadratic trace identities}\label{app:quadratic}
Let $\chi$ be the quadratic character of $\F_p$, extended by zero, and put $T_j(p)=\sum_{\lambda\ne0,1}\chi(\lambda)a_p(\lambda)^j$. The following is a known-trace-formula corollary.
\begin{theorem}\label{thm:quadratic}
For every prime $p\ge5$,
\[
 T_4(p)=-c(p)-b(p)-3p\,a(p)-2p^2-3p-1,
\]
where $a=8.4.a.a=\eta(2z)^4\eta(4z)^4$,
$b=4.6.a.a=\eta(2z)^{12}$, and
$c=8.6.a.a=q+20q^3-74q^5-24q^7+\cdots$.
\end{theorem}
\begin{proof}
Use the Ahlgren and Ahlgren--Ono trace formulas as stated in
\cite[Theorem 2.3]{Grove}:
\begin{align*}
 -\Tr_{k+2}(\Gamma_0(4),p)&=3+\sum_{\lambda\ne0,1}G_{k+2}(a(\lambda),p),\\
 -\Tr_{k+2}(\Gamma_0(8),p)&=4+\sum_{\gamma\ne0,1,-1}G_{k+2}(a(\gamma^2),p).
\end{align*}
Here $G_4(s,p)=s^2-p$ and $G_6(s,p)=s^4-3ps^2+p^2$.
Put $U_k=\sum_{\lambda\ne0,1}a(\lambda)^k$ and
$Q_k=\sum_{\gamma\ne0,1,-1}a(\gamma^2)^k$. Then
$Q_k-U_k=T_k$, and $Q_0-U_0=-1$.
The relevant spaces have bases
\[
 S_4(\Gamma_0(4))=0,\quad S_4(\Gamma_0(8))=\langle a\rangle,
\quad S_6(\Gamma_0(4))=\langle b\rangle,
\]
\[
 S_6(\Gamma_0(8))=\langle b(z),b(2z),c(z)\rangle.
\]
Subtracting the weight-four trace formulas gives
$-a(p)=1+T_2(p)+p$, so $T_2=-a(p)-p-1$.
Subtracting the weight-six formulas gives
$-(b(p)+c(p))=1+T_4-3pT_2-p^2$.
Substitution proves the formula.
\end{proof}

The identity is a consequence of established trace formulas. No new
modularity or priority assertion is attached to this appendix.

\section{The corrected hypergeometric normalization}\label{sec:hyp}
Assume $p\ge5$, $p\equiv1\pmod3$. Let $\phi$ and $\rho$ be the quadratic and a primitive cubic character, and let $g$ and $J$ denote Gauss and Jacobi sums for a fixed nontrivial additive character. We use the normalized Gauss-sum definition
\cite[Eq. (1.1)]{Grove}, with all multiplicative characters, including
the trivial character, extended by zero at zero. Put
\[
 H(\lambda)=\frac1{1-p}\sum_\chi c(\chi)\chi(\lambda),
 \qquad c(\chi)=\frac{g(\phi\chi)^2g(\bar\chi)^2}{g(\phi)^2}.
\]
Only nonzero arguments are used here; the separately defined value at zero is immaterial. For nonsingular $\lambda$, $a(\lambda)=\phi(-1)H(\lambda)$.
Fix a generator $\omega$ of the multiplicative character group such that $\rho=\omega^{(p-1)/3}$ and $\phi=\omega^{(p-1)/2}$. The same $\omega$ is used to interpret the rational parameters of $H_p(A;1)$. Let
\[
 A=\left\{\tfrac12,\tfrac12,\tfrac13,\tfrac13;
               1,1,\tfrac56,\tfrac56\right\},\qquad
 J=J(\rho,\phi\bar\rho).
\]
\begin{lemma}\label{lem:norm}
With this normalization,
\begin{equation}\label{eq:correctnorm}
 M_{3,2}(p)+1
 =\frac1{p-1}\sum_\chi c(\chi)c(\bar\rho\bar\chi)
 =-J^2 H_p(A;1).
\end{equation}
\end{lemma}
\begin{proof}
Since $\phi$ is nontrivial,
$c(\chi)=J(\phi\chi,\bar\chi)^2$ by the Gauss--Jacobi identity.
Expanding the Jacobi sums and summing over $\chi$ imposes
$xy=(1-x)(1-y)$, equivalently $x+y=1$. Therefore
\[
 \sum_\chi c(\chi)=(p-1)\sum_x\phi(x(1-x))
 =-(p-1)\phi(-1),
\]
so $H(1)=\phi(-1)$. Character orthogonality applied to $H(\lambda)^2$
then gives the first equality of \eqref{eq:correctnorm}.

The numerator of each summand for $H_p(A;1)$ is
\[
 g(\phi\chi)^2g(\rho\chi)^2g(\bar\chi)^2
                  g(\phi\bar\rho\bar\chi)^2.
\]
Its normalizing denominator is
$g(\phi)^2g(\rho)^2g(\varepsilon)^2g(\phi\bar\rho)^2$,
where $g(\varepsilon)^2=1$. The convolution on the middle of
\eqref{eq:correctnorm} has denominator $g(\phi)^4$ and the opposite
overall sign, since $1-p=-(p-1)$. Their ratio is
\[
 -\left(\frac{g(\rho)g(\phi\bar\rho)}{g(\phi)}\right)^2=-J^2.
\]
This proves the normalization factor.
\end{proof}

For instance, choose $\omega(3)=e^{2\pi i/6}$ at $p=7$ and
$\rho=\omega^2$. Then $J=2+\sqrt{-3}$, $J^2=1+4\sqrt{-3}$,
and $M_{3,2}(7)+1=-15$. Hence
$H_7(A;1)=15/(1+4\sqrt{-3})$, which is not $-15$.
The v1 phrase identifying the convolution directly with the normalized
$A$-value must therefore include the factor in \eqref{eq:correctnorm}.
The underlying moment calculation is unchanged.

Under the carried-certificate comparison of Section~\ref{sec:cubic}, the corrected modular relation is
\[
 J^2 H_p(A;1)=a_f(p)+p.
\]
For $B=\{\frac12,\frac12,\frac13,\frac23;1,1,1,1\}$, the
published identity $H_p(B;1)=a_f(p)+\leg{3}{p}p$
\cite[Theorem~2 and Table~1]{LTYZ} gives
\[
 J^2 H_p(A;1)=H_p(B;1)+\left(1-\leg{3}{p}\right)p.
\]
The datum $A$ also occurs in the EHMM literature at a different
argument \cite{EHMM2}. A priority claim requires comparison of both
the argument and the Jacobi twists. The algebra above supplies the normalization proof; numerical scripts are retained separately as historical checks.


\section{An analytic period substitution and the row-sign correction}\label{app:period}

This appendix records a change of variables and a sign convention. Relations explicitly marked numerical refer to the historical candidate model and are not critical-value theorems.

\subsection{An exact change of variables}
Let
\[
 F(\lambda)={}_2F_1\!\left(\tfrac12,\tfrac12;1;\lambda\right),
 \qquad F_c(\lambda)=F(1-\lambda),\qquad 0<\lambda<1.
\]
For an integer $k\geq0$, an integer $0\leq j\leq k$, and a real number
$e>0$, define $y=F_c/F$. The inverse parameter is the modular lambda
function at $iy$. Put
\[
 G_e(iy)=\lambda(iy)^e\bigl(1-\lambda(iy)\bigr)
          \theta_3(0\mid iy)^{2k+4}.
\]
Then the following analytic identity holds:
\begin{equation}\label{eq:period-dictionary-static}
 \int_0^1\lambda^{e-1} F(\lambda)^j F_c(\lambda)^{k-j}\,d\lambda
 =\pi\int_0^\infty G_e(iy)y^{k-j}\,dy.
\end{equation}

\begin{proof}
Both $u=F$ and $v=F_c$ solve
\[
 \lambda(1-\lambda)u''+(1-2\lambda)u'-\tfrac14u=0.
\]
This is the specialization of the hypergeometric differential equation in
DLMF~15.10.1 \cite{DLMF}. Abel's identity therefore gives
\[
 uv'-u'v=\frac{C}{\lambda(1-\lambda)}.
\]
The hypergeometric--elliptic identity and the logarithmic endpoint expansion
(DLMF~19.5.1 and 19.12.1) yield
\[
 F(\lambda)=1+O(\lambda),\qquad
 F_c(\lambda)=\frac1\pi\log\frac{16}{\lambda}
                +O(\lambda|\log\lambda|)
 \quad(\lambda\longrightarrow0^+).
\]
The differentiated expansions imply $C=-1/\pi$. Hence
\begin{equation}\label{eq:period-jacobian-static}
 \frac{dy}{d\lambda}
 =-\frac{1}{\pi\lambda(1-\lambda)F(\lambda)^2}.
\end{equation}
The positive function $y$ decreases from $+\infty$ to $0$. DLMF~20.9.2,
together with DLMF~19.5.1, identifies
$F(\lambda(iy))=\theta_3(0\mid iy)^2$ and
$F_c(\lambda(iy))/F(\lambda(iy))=y$.
Using $F_c=yF$ and
$d\lambda=-\pi\lambda(1-\lambda)F^2\,dy$ proves
\eqref{eq:period-dictionary-static} after reversing the endpoints.
The original integrand is
$O(\lambda^{e-1}(1+|\log\lambda|)^{k-j})$ near $0$ and
$O((1+|\log(1-\lambda)|)^j)$ near $1$.
Both bounds are integrable under the stated hypotheses, so the improper
change of variables is justified.\end{proof}

For the present candidate periods, $k=4$ and
$e\in\{1/3,2/3,4/3,5/3\}$; writing $m=k+1-j$ gives the factor
$y^{m-1}$ used in the old note N12. The same argument also covers its
$k=2$ control integrals.
This establishes the analytic dictionary and its normalization. Identifying
its right-hand side with a specified arithmetic $L$-value requires further
modularity, rational-structure, and Mellin-normalization arguments. In
particular, the rational Scholl realization in Section~\ref{sec:scholl}
does not by itself identify these integrals with rational period bases or
prove the rational determinant constants below.

\subsection{The documented and implemented signs differ}
Write $H_j(\lambda)=F(\lambda)^jF_c(\lambda)^{4-j}$ and
\[
 C_c(j)=\int_0^1(\lambda^c+\lambda^{-c})H_j(\lambda)\,d\lambda.
\]
For two row families $A_{2/3},A_{1/3}$ use the fixed row order
\[
 \mu_A(j_1,j_2)=
 \det\begin{pmatrix}
 A_{2/3}(j_1)&A_{2/3}(j_2)\\
 A_{1/3}(j_1)&A_{1/3}(j_2)
 \end{pmatrix}.
\]
N12, line~45, describes the uniform convention
\[
 S_c^{\mathrm{unif}}(j)
 =\int_0^1(\lambda^c-\lambda^{-c})H_j(\lambda)\,d\lambda.
\]
The actual N11 implementation instead uses
\begin{align*}
 S_{2/3}^{\mathrm{code}}(j)
   &=\int_0^1(\lambda^{-2/3}-\lambda^{2/3})H_j(\lambda)\,d\lambda,\\
 S_{1/3}^{\mathrm{code}}(j)
   &=\int_0^1(\lambda^{1/3}-\lambda^{-1/3})H_j(\lambda)\,d\lambda.
\end{align*}
Thus, exactly and without a numerical computation,
\[
 \mu_{S^{\mathrm{code}}}(j_1,j_2)
 =-\mu_{S^{\mathrm{unif}}}(j_1,j_2).
\]
The old negative ratios use the implemented convention. To retain the
uniform convention in prose, the corresponding \emph{reported numerical}
relations must instead read
\begin{align*}
 \frac{\mu_{S^{\mathrm{unif}}}(1,3)}{\mu_C(2,4)}
     &\mathrel{\stackrel{\mathrm{num}}{=}}\frac12, &
 \frac{\mu_{S^{\mathrm{unif}}}(0,2)}{\mu_C(1,3)}
     &\mathrel{\stackrel{\mathrm{num}}{=}}\frac43,\\
 \mathcal L_{\rm candidate}(4)&\mathrel{\stackrel{\mathrm{num}}{=}}
       \frac{\pi^6\mu_{S^{\mathrm{unif}}}(0,2)}{209952}.
\end{align*}
Alternatively, explicitly define the two implemented row orientations and
retain the negative signs in the historical outputs. No numerical rerun is
needed to diagnose this discrepancy. The five principal critical-value
relations expressed only through $\mu_C$ are unaffected.


\section*{Revision record and computational supplement}
The public predecessors are the original preprint \cite{WalkeV1} and
its version-2 consolidation \cite{WalkeV2}. This revision integrates the
local arguments, conductor deductions and irreducibility argument into
one manuscript. The earlier versions remain historical records; their
identifiers do not identify the text of this revision.

The earlier corrections repaired the quartic finite test set,
hypergeometric normalization and numerical validators. The cubic proof
supplied the geometric inputs for character order three. The good-prime
purity and rational full-level-12 Scholl construction supplied the weight,
projectors and Hodge numbers. This revision retains those arguments and adds:
\begin{itemize}
\item the explicit local-at-$3$ quotient, inertia calculation and
Frobenius sign for the literal plus descent;
\item the true arithmetic Galois closure at $2$, its semistable model,
a full tameness argument, the tail vanishing and the equivariant
nearby-cycle calculation;
\item the resulting conductor products and their auxiliary-prime scope;
\item strong absolute irreducibility for $\ell\ne3$, with its precise
local and Hodge--Tate inputs;
\item the current claim/dependency map in Section~\ref{sec:claims}.
\end{itemize}
The true dyadic splitting cover is $\eta^3=4(v^3+1)$. The earlier
non-Galois cover $s^3=v^3-1$ was a valid tameness device; it was not the
arithmetic Galois closure needed for the later inertia calculation.
Here the complete true-cover charts also meet the tameness lemma's
hypotheses, so the local proof uses that one cover throughout.

The package contains the current TeX sources, current claim/status files,
small exact local-check scripts with their retained run reports, and the
unchanged earlier source/supplement archive. That historical archive
contains the bounded cubic certificate and its independent exact
implementation, along with older numerical records. Its earlier open-task
lists and metadata describe their original stage and are superseded by
the current claim map, not silently rewritten. Source hashes establish
file preservation, not mathematical truth.
The finite algebra checks do not implement nearby cycles or certify the
geometric proof. The numerical $L$-function experiments and large scans
have not been rerun for this consolidation. Both earlier success and
failure records remain historical evidence.

\subsection*{AI assistance and review status}
The earlier project used ChatGPT and Claude, as disclosed in v1.
AI assistance during preparation and revision extended throughout the
manuscript, including the abstract, exposition, mathematical arguments
and appendices. Its uses included drafting and revising mathematical
arguments, preparing code and algebraic and sign checks, looking up
references, and preparing and revising the LaTeX source. In particular,
the new local, conductor and irreducibility arguments in
Sections~\ref{sec:local-three}--\ref{sec:irreducibility} were drafted and
internally checked with ChatGPT assistance. This internal AI-assisted
checking does not constitute independent specialist review. Model
versions are not specified because they are not verified in the
available preparation record. AI systems are not authors.
Responsibility for the arguments, references, computational record
and disclosure rests with the author. No external specialist review,
journal acceptance, proof-assistant verification, or exhaustive novelty
determination is represented by this manuscript.
The newly written local and irreducibility proofs do not establish
rank-four automorphy, analytic continuation, the functional-equation sign,
arithmetic period identities, or a classical elliptic BSD theorem.

\begin{thebibliography}{99}
\bibitem{WalkeV1} R. M. Walke,
\emph{Twisted moments of the Legendre family, rigid Calabi--Yau modular forms,
and a weight-five motive of conductor 1296}, Zenodo preprint v1,
28 October 2026. \url{https://doi.org/10.5281/zenodo.23719224}.
This is the historical public predecessor.
\bibitem{SGA1} A. Grothendieck et al.,
\emph{Revetements etales et groupe fondamental (SGA1)}, Expose XIII,
Definition 2.1.1 and Corollary 2.9, recomposed edition.
\url{https://arxiv.org/pdf/math/0206203}.
\bibitem{SGA4} M. Artin, A. Grothendieck and J.-L. Verdier (eds.),
\emph{Theorie des topos et cohomologie etale des schemas (SGA4)}, tome 3,
Expose XVIII, Theorem 2.9, Var. 2, and Theorem 3.2.5.
\url{https://www.cmls.polytechnique.fr/perso/laszlo/sga4/SGA4-3/sga43.pdf}.
\bibitem{Katz} N. M. Katz,
\emph{Convolution and Equidistribution: Sato--Tate Theorems for
Finite-Field Mellin Transforms}, Chapter 15.
\url{https://web.math.princeton.edu/~nmk/mellin331.pdf}.
\bibitem{Tate} J. Tate, \emph{A review of non-Archimedean elliptic functions},
section on points of finite order, pp. 18--19 of the manuscript.
\url{https://people.math.harvard.edu/~ctm/home/text/others/tate/tate-unif/tate-unif.pdf}.
\bibitem{Milne} J. S. Milne, \emph{Lectures on Etale Cohomology},
especially Theorem 24.1 and the sections on adic cohomology.
\url{https://www.jmilne.org/math/CourseNotes/LECc.pdf}.
\bibitem{StacksTrace} The Stacks Project, Theorem 64.20.2, tag 03V0,
finite-field adic trace formula. \url{https://stacks.math.columbia.edu/tag/03V0}.
\bibitem{StacksFinite} The Stacks Project, Proposition 59.55.2, tag 03QP,
vanishing of finite higher direct images.
\url{https://stacks.math.columbia.edu/tag/03QP}.
\bibitem{StacksDescent} The Stacks Project, Section 59.104, tag 0GEX,
descent for etale sheaves. \url{https://stacks.math.columbia.edu/tag/0GEX}.
\bibitem{DGP} L. Dieulefait, L. Guerberoff and A. Pacetti,
\emph{Proving modularity for a given elliptic curve over an imaginary
quadratic field}, Theorem 4.1. \url{https://arxiv.org/abs/0804.2302}.
\bibitem{Sturm} J. Sturm, \emph{On the congruence of modular forms},
Lecture Notes in Math. \textbf{1240}, Springer, 1987, 275--280,
Theorem~1. \url{https://doi.org/10.1007/BFb0072985}.
\bibitem{LTYZ} L. Long, F.-T. Tu, N. Yui and W. Zudilin,
\emph{Supercongruences for rigid hypergeometric Calabi--Yau threefolds},
Adv. Math. \textbf{393} (2021), 108058.
\url{https://arxiv.org/abs/1705.01663}.
\bibitem{KS} M. Kazalicki and A. J. Scholl,
\emph{Modular forms, de Rham cohomology and congruences},
Trans. Amer. Math. Soc. \textbf{368} (2016), 7097--7117, Sections 7--8.
\url{https://www.dpmms.cam.ac.uk/~ajs1005/preprints/mfdrc.pdf}.
\bibitem{Scholl} A. J. Scholl,
\emph{On some $\ell$-adic representations of
$\operatorname{Gal}(\overline\Q/\Q)$ attached to noncongruence subgroups},
Bull. London Math. Soc. \textbf{38} (2006), 561--567.
\url{https://www.dpmms.cam.ac.uk/~ajs1005/preprints/galois-2004.pdf}.
\bibitem{ALLL} A. O. L. Atkin, W.-C. W. Li, T. Liu and L. Long,
\emph{Galois representations with quaternion multiplication associated
to noncongruence modular forms}, Trans. Amer. Math. Soc.
\textbf{365} (2013), 6217--6242.
\url{https://www.math.purdue.edu/~tongliu/pub/tran6019.pdf}.
\bibitem{Grove} B. Grove, \emph{Hypergeometric moments and Hecke trace formulas},
Ramanujan J. \textbf{69} (2026), article 99.
\url{https://doi.org/10.1007/s11139-026-01372-y};
\url{https://arxiv.org/abs/2409.14502}.
\bibitem{AhlgrenOno} S. Ahlgren and K. Ono,
\emph{Modularity of a certain Calabi--Yau threefold},
Monatsh. Math. \textbf{129} (2000), 177--190.
\url{https://uva.theopenscholar.com/files/ken-ono/files/050.pdf}.
\bibitem{Ahlgren} S. Ahlgren,
\emph{The points of a certain fivefold over finite fields and the twelfth
power of the eta function}, Finite Fields Appl. \textbf{8} (2002), 18--33.
\bibitem{EHMM2} M. Allen, B. Grove, L. Long and F.-T. Tu,
\emph{The explicit hypergeometric-modularity method II},
Res. Math. Sci. \textbf{12} (2025), article~84.
\url{https://doi.org/10.1007/s40687-025-00566-y}.
Preprint: \url{https://arxiv.org/abs/2411.15116}.
\bibitem{DLMF} NIST Digital Library of Mathematical Functions,
Equations 15.10.1, 19.5.1, 19.12.1 and 20.9.2.
\url{https://dlmf.nist.gov/15.10.E1},
\url{https://dlmf.nist.gov/19.5.E1},
\url{https://dlmf.nist.gov/19.12.E1},
\url{https://dlmf.nist.gov/20.9.E2}.

\bibitem{DeligneWII} P. Deligne, \emph{La conjecture de Weil: II},
Publ. Math. IHES \textbf{52} (1980), 137--252,
Corollary~3.3.6 (p.~206); see also Section~3.7.1.
\url{https://publications.ias.edu/sites/default/files/Number40.pdf}.
\bibitem{SchollMotives} A. J. Scholl, \emph{Motives for modular forms},
Invent. Math. \textbf{100} (1990), 419--430.
\url{https://www.dpmms.cam.ac.uk/~ajs1005/preprints/mf.pdf}.
Author's corrections: \url{https://www.dpmms.cam.ac.uk/~ajs1005/preprints/mf_corr.html}.
\bibitem{SchollVanishing} A. J. Scholl,
\emph{Vanishing cycles and non-classical parabolic cohomology},
Invent. Math. \textbf{124} (1996), 503--524.
\url{https://www.dpmms.cam.ac.uk/~ajs1005/preprints/van.pdf}.
Correction: \url{https://www.dpmms.cam.ac.uk/~ajs1005/preprints/van-corr/van-corr.pdf}.
\bibitem{StacksMinimal} The Stacks Project, Lemma 55.10.1, tag 0C6B,
uniqueness of the minimal model.
\url{https://stacks.math.columbia.edu/tag/0C6B}.
\bibitem{AntieauMeier} B. Antieau and L. Meier,
\emph{The Brauer group of the moduli stack of elliptic curves},
Algebra Number Theory \textbf{14} (2020), 2295--2333,
Lemma~4.3 and Proposition~4.5.
\url{https://arxiv.org/pdf/1608.00851}.
\bibitem{WalkeV2} R. M. Walke, \emph{Cubic twists of Legendre moments:
cohomological descent and Frobenius trace formulas}, version 2 (2026).
\url{https://doi.org/10.5281/zenodo.23775155}.
This is the public normalization predecessor of the present research draft.

\bibitem{Ulmer} D. Ulmer, \emph{Conductors of $\ell$-adic representations},
Proc. Amer. Math. Soc. 144 (2016), 2291--2299, \S\S5--9.
\url{https://arxiv.org/abs/1307.4525}.

\bibitem{Nakayama} C. Nakayama, \emph{Nearby cycles for log smooth families},
Compositio Mathematica \textbf{112} (1998), 45--75,
Variant 3.4.1, Proposition 4.2 and Remark 4.2.1.
\url{https://doi.org/10.1023/A:1000327225021}.

\bibitem{SerreRankOne} J.-P. Serre, \emph{Abelian $\ell$-adic Representations and
Elliptic Curves}, 1968; Chapter III, \S1.2, Theorem (Tate), printed p. III-7,
and Appendix. Primary scan:
\url{https://www.math.mcgill.ca/darmon/courses/18-19/gs/serre-mcgill.pdf}.

\bibitem{Yun} Z. Yun, with an appendix by C. Vincent,
\emph{Galois representations attached to moments of Kloosterman sums and
conjectures of Evans}, Compositio Math. \textbf{151} (2015), 68--120,
Sections 2.5.5 and 4.3.1--4.3.4.
\url{https://arxiv.org/abs/1308.3920}.

\bibitem{Haessig} C. D. Haessig,
\emph{Kloosterman sums and Hecke polynomials in characteristics 2 and 3},
Finite Fields Appl. \textbf{76} (2021), 101896, proof of Theorem 2.2.
\url{https://arxiv.org/abs/2012.00570}.

\bibitem{SGA7} P. Deligne and N. Katz (eds.),
\emph{Groupes de monodromie en geometrie algebrique (SGA 7 II)},
Lecture Notes in Math. \textbf{340}, Springer (1973),
Expose XIII, Lemmas 2.1.10--2.1.11.
\url{https://library.slmath.org/nonmsri/sga/sga/pdf/sga7-2.pdf}.

\bibitem{Illusie} L. Illusie, \emph{An overview of the work of K. Fujiwara,
K. Kato, and C. Nakayama on logarithmic etale cohomology}, Asterisque
\textbf{279} (2002), 271--322, Corollaries 8.4.3--8.4.4.
\url{https://www.numdam.org/article/AST_2002__279__271_0.pdf}.
\end{thebibliography}

\end{document}
